%% This document created by Scientific Word (R) Version 3.0


\documentclass[11pt]{article}
\usepackage{amssymb}

%%%%%%%%%%%%%%%%%%%%%%%%%%%%%%%%%%%%%%%%%%%%%%%%%%%%%%%%%%%%%%%%%%%%%%%%%%%%%%%%%%%%%%%%%%%%%%%%%%%%
\usepackage{graphicx}
\usepackage{amsmath}

%TCIDATA{OutputFilter=LATEX.DLL}
%TCIDATA{Created=Wed Oct 27 14:48:41 1999}
%TCIDATA{LastRevised=Fri Jun 09 09:56:05 2000}
%TCIDATA{<META NAME="GraphicsSave" CONTENT="32">}
%TCIDATA{<META NAME="DocumentShell" CONTENT="General\Blank Document">}
%TCIDATA{Language=American English}
%TCIDATA{CSTFile=LaTeX article (bright).cst}
%TCIDATA{PageSetup=57,43,0,0,0}
%TCIDATA{Counters=arabic,1}
%TCIDATA{AllPages=
%H=36
%F=36
%}


\newtheorem{theorem}{Theorem}
\newtheorem{acknowledgement}[theorem]{Acknowledgement}
\newtheorem{algorithm}[theorem]{Algorithm}
\newtheorem{axiom}[theorem]{Axiom}
\newtheorem{case}[theorem]{Case}
\newtheorem{claim}[theorem]{Claim}
\newtheorem{conclusion}[theorem]{Conclusion}
\newtheorem{condition}[theorem]{Condition}
\newtheorem{conjecture}[theorem]{Conjecture}
\newtheorem{corollary}[theorem]{Corollary}
\newtheorem{criterion}[theorem]{Criterion}
\newtheorem{definition}[theorem]{Definition}
\newtheorem{example}[theorem]{Example}
\newtheorem{exercise}[theorem]{Exercise}
\newtheorem{lemma}[theorem]{Lemma}
\newtheorem{notation}[theorem]{Notation}
\newtheorem{problem}[theorem]{Problem}
\newtheorem{proposition}[theorem]{Proposition}
\newtheorem{remark}[theorem]{Remark}
\newtheorem{solution}[theorem]{Solution}
\newtheorem{summary}[theorem]{Summary}
\newenvironment{proof}[1][Proof]{\textbf{#1.} }{\ \rule{0.5em}{0.5em}}
\input{tcilatex}

\begin{document}

\title{Boundary Conditions in the B2.5 Code}
\author{Mladen Stanojevi\'{c}}
\maketitle
\tableofcontents

\section{Introduction}

LUD 24 Dec 99\bigskip

The boundary conditions are implemented in the B2.5 code by means of the
boundary sources, that is in a similar way as the volume sources due to the
atomic processes, but they are defined only in the boundary cells of the
computational grid. [\ref{mbaelmans}] It means that in a case of the
boundary conditions at the material surfaces the boundary cells have a role
of combined electrostatic sheath and quasineutral collisionless magnetic
presheath, which cannot be modelled with the B2.5 fluid equations. According
to the existing fluid models of the magnetized boundary plasma, the role of
the magnetic presheath, which is located between the collisional plasma
presheath and electrostatic sheath, is to deflect the charged particles from
their average motion along the magnetic field lines to the direction
perpendicular to the wall. The average ion velocity along the magnetic field
direction at the entrance of the magnetic presheath from the collisional
presheath side is at least the local ion sound velocity, that is $c_{sMPSE}=%
\sqrt{\frac {k_{B}(ZT_{eMPSE}+T_{iMPSE})}{m_{i}}}$ in a two-component
plasma, or $u_{\Vert MPSE}\geq c_{sMPSE}$ [\ref{mbohm95}]. The component of
the average ion velocity at the entrance of the electrostatic sheath, which
is perpendicular to the wall, is also greater than or equal to the local ion
sound velocity, i.e., $u_{SE}\geq c_{sSE}$. [\ref{mriemann91}, \ref{mbohm95}%
]\ Since both the magnetic presheath and electrostatic sheath are
collisionless regions and assuming that the dissipative effects due to some
turbulent mechanisms are negligible in these regions, particle, momentum and
energy fluxes are conserved in these regions. The B2.5 equations for the
bulk plasma are approximately valid in the collisional presheath, assuming
that the charged particle velocity distributions do not differ significantly
from Maxwellian or drifting-Maxwellian distributions. (However, it should be
emphasized that according to the two-scale kinetic analysis of the boundary
plasma near the solid surface, pure Maxwellian or drifting-Maxwellian ion
velocity distribution cannot satisfy the generalized Bohm criterion. [\ref
{mriemann91}, \ref{mriemann97}, \ref{mdaube99}]) It means that they provide
the values for the charged particle densities, drift velocities, ion and
electron temperatures at the magnetic presheath edge ($MPSE$), which can be
used for calculating the fluxes at the magnetic presheath edge. The value of
the poloidal component of the ion drift velocity at the magnetic presheath
edge obtained from the bulk transport equations should be corrected, because
of the effects related to the stronger local electric field in the magnetic
presheath. This electric field causes a strong acceleration of the ions
along the magnetic field lines, similar like in a case of a non-magnetic
plasma presheath or in a case when the magnetic field is perpendicular to
the surface, and additional poloidal $\mathbf{E\times B}$ drift due to the
radial component of the electric field. In some cases $\mathbf{E\times B}$
drift velocity can be comparable with the local ion sound velocity. [\ref
{mstangeby95}, \ref{mrognlien1}, \ref{mrognlien2}] The correction of the
poloidal component of the ion drift velocity due to the $\mathbf{B}%
\times\nabla p$ drift is usually negligible, since it contributes mostly to
the ion flow parallel to the surface. [\ref{mchankin94}] However, more
detailed investigations of the importance of this drift should be performed
in the future studies, which should include also kinetic models and
numerical simulations (e.g. particle and/or kinetic simulations) of the
magnetized boundary plasma.

\bigskip

The boundary sources in the B2.5 code are defined by the boundary cell
averages or volume integrals of the negative divergences of the incoming
fluxes into the boundary cells. The boundary condition for the continuity
equation for each atomic species is given by: 
\begin{equation}
S_{p}=-\int_{B.V.}\nabla\cdot\mathbf{\Gamma}\text{ }dV=-\int_{B.S.}\mathbf{%
\Gamma\cdot} d\mathbf{A}\text{,}   \label{1}
\end{equation}
which means that in the discretized form of this equation the value of the
particle flux, $\Gamma_{j}=n\left\langle v_{j}\right\rangle $, at the
magnetic presheath edge as a function of the macroscopic plasma parameters
should be specified. The boundary condition for the momentum equation is: 
\begin{equation}
\mathbf{S}_{m}=-\int_{B.V.}\nabla\cdot\underline{\mathbf{M}}\text{ }%
dV=-\int_{B.S.}\underline{\mathbf{M}}\cdot d\mathbf{A}\text{.}   \label{2}
\end{equation}
In this case the appropriate components of the momentum tensor, $%
M_{jk}=mn\left\langle v_{j}v_{k}\right\rangle $, at the magnetic presheath
edge should be specified in the discretized equations. Similarly, the energy
or heat boundary sources are given by: 
\begin{equation}
S_{q}=-\int_{B.V.}\nabla\cdot\mathbf{Q}\text{ }dV=-\int_{B.S.}\mathbf{Q\cdot}%
d\mathbf{A}\text{,}   \label{3}
\end{equation}
where the appropriate components of the electron and/or ion heat flux, $%
Q_{j}=\frac{1}{2}mn\left\langle v^{2}v_{j}\right\rangle $, at the magnetic
presheath edge as functions of the macroscopic plasma parameters should be
specified in the discretized equations. The boundary charge source is given
by: 
\begin{equation}
S_{ch}=-\int_{B.V.}\nabla\cdot\mathbf{j}\text{ }dV=-\int_{B.S.}\mathbf{j\cdot%
}d\mathbf{A}\text{,}   \label{4}
\end{equation}
where the electric current density, $\mathbf{j}=e\left( \sum_{\alpha
}Z_{\alpha}\mathbf{\Gamma}_{\alpha}-\mathbf{\Gamma}_{e}\right) $, at the
magnetic presheath edge should be specified in the discretized equations.

\bigskip

\section{Analytical Expressions for the Standard (``Old'') Boundary Sources
in the B2.5 Code}

\bigskip

The boundary sources are expressed as linear and/or weakly non-linear
functions of the macroscopic plasma parameters ($n_{\alpha}$, $n_{e}$, $%
u_{\alpha x}$, $\widehat{T}_{i}$, $\widehat{T}_{e}$, $\Phi$), calculated in
the cells neighbouring to the boundary cells. Specific dependences of the
boundary sources on these parameters are determined by the choice of the
coefficients $C_{pj\alpha}$, $C_{mj\alpha}$, $C_{hij\alpha}$, $C_{hej}$ and $%
C_{chj}$ (see Section 3). These coefficients are free input parameters,
which depend on the theoretical model used for describing the boundary
conditions. They must be written by the user in the input file \texttt{%
b2ah.dat}. If \texttt{b2ah.dat} file already exists and the boundary
conditions for the new simulation are not significantly different, it is
practical to write the minor changes of the boundary conditions in the input
file \texttt{b2mn.dat}. (If the existing \texttt{b2ah.dat} file is
corrected, it is not necessary to do that.) When the theoretical model
predicts more complicated non-linear dependence of the boundary fluxes on
the plasma parameters, the effects of these non-linearities can be only
approximately taken into consideration in the calculation, e.g., by some
appropriate choice of the numerical values of the input parameters (see the
example in Section 3). More details about the numerical implementation of
the boundary sources can be found in the \texttt{b2stbc.f} source code.

\bigskip

\subsection{Definitions of Physical Parameters Used in the Analytical
Expressions for the Standard Boundary Sources}

\bigskip

All following quantities are calculated locally at the appropriate boundary
- SOUTH (S), NORTH (N) , WEST (W) and EAST (E) (see Fig. \ref{fig1}, AB and
CD are the divertor plates, $x$ is the poloidal coordinate, $y$ is the
radial coordinate, the separatrix goes from W through F and F$^{%
%TCIMACRO{\UNICODE{0xb4}}%
%BeginExpansion
{\acute{}}%
%EndExpansion
}$ to E [\ref{mbaelmans}]), i.e., in the boundary cells:

\bigskip

\begin{equation}
S_{0}=V/h_{y1}\text{, \ \ \ }h_{y1}=h_{y}\cos(\widehat{x}\widehat{y}) 
\label{5}
\end{equation}
for the SOUTH and NORTH boundary, where $V$ is the volume of the boundary
cell, $h_{y}$is the y-diameter (radial length) of the boundary cell and $%
\cos(\widehat{x}\widehat{y})$ is the cosine of the complementary angle
between the x-direction and y-direction on the boundary cell;

\bigskip

\begin{equation}
S_{0}=V/h_{x}   \label{6}
\end{equation}
for the WEST and EAST boundary, where $V$ is the volume of the boundary cell
and $h_{x}$ is the x-diameter (poloidal length) of the boundary cell;

\bigskip

\begin{equation}
S_{1}=\left| S_{1}^{\pm}\right| =\left| \frac{B_{y}}{B}A\right|   \label{7}
\end{equation}
for the SOUTH and NORTH boundary, where $S_{1}^{\pm}$ is the product of the
magnetic field pitch and the surface area on the cell face between the
boundary cell and its neighbouring (previous) cell ((ix,-1) and (ix,0) for
the SOUTH and (ix,ny) and (ix,ny-1) for the NORTH);

\bigskip

\begin{equation}
S_{1}=\left| S_{1}^{\pm}\right| =\left| \frac{B_{x}}{B}A\right|   \label{8}
\end{equation}
for the WEST and EAST boundary, where $S_{1}^{\pm}$ is the product of the
magnetic field pitch and the surface area on the cell face between the
boundary cell and its neighbouring (previous) cell ((-1,iy) and (0,iy) for
the WEST and (nx,iy) and (nx-1,iy) for the EAST);

\bigskip

\begin{equation}
c_{s\alpha}=\sqrt{\frac{k_{B}\left( Z_{\alpha}T_{e}+T_{i}\right) }{%
A_{\alpha}m_{p}}}=\sqrt{\frac{k_{B}\left( Z_{\alpha}T_{e}+T_{i}\right) }{%
m_{\alpha}}}   \label{9}
\end{equation}
is the isothermal ion sound velocity for the atomic species $\alpha$;

\bigskip

\begin{equation}
c_{st}=\sqrt{\frac{\sum_{\alpha}n_{\alpha}k_{B}\left(
Z_{\alpha}T_{e}+T_{i}\right) }{\sum_{\alpha}n_{\alpha}A_{\alpha}m_{p}}}=%
\sqrt{\frac {\sum_{\alpha}n_{\alpha}k_{B}\left( Z_{\alpha}T_{e}+T_{i}\right)
}{\sum_{\alpha}n_{\alpha}m_{\alpha}}}   \label{10}
\end{equation}
is the isothermal ion sound velocity for multispecies plasma;

\bigskip

\begin{equation}
c_{sD\alpha}=\max(c_{s\alpha},0.1c_{s\alpha}\pm\left( B_{z}/B_{x}\right)
u_{Dx})   \label{11}
\end{equation}
is conditionally corrected ion drift velocity of the atomic species $\alpha$
due to the poloidal $\mathbf{E\times B}$ drift, where ''$+$'' is for the
WEST and ''$-$'' is for the EAST; the approximate correction is used only if 
$0.1c_{s\alpha}\pm\left( B_{z}/B_{x}\right) u_{Dx}>c_{s\alpha}$;

\bigskip

\begin{equation}
c_{stD}=\max(c_{st},0.1c_{st}\pm\left( B_{z}/B_{x}\right) u_{Dx}) 
\label{12}
\end{equation}
is conditionally corrected ion drift velocity of the multispecies plasma due
to the poloidal $\mathbf{E\times B}$ drift, where ''$+$'' is for the WEST
and ''$-$'' is for the EAST; the approximate correction is used only if $%
0.1c_{st}\pm\left( B_{z}/B_{x}\right) u_{Dx}>c_{st}$;

\bigskip

\begin{equation}
\rho_{\alpha}=A_{\alpha}m_{p}n_{\alpha}=m_{\alpha}n_{\alpha}   \label{13}
\end{equation}
is the mass density of the atomic species $\alpha$;

\bigskip

\begin{equation}
n_{i}=\sum_{\alpha}n_{\alpha}   \label{14}
\end{equation}
is the total atomic species (ion) particle density;

\bigskip

\begin{equation}
t_{\alpha}=\frac{n_{\alpha}}{n_{i}}   \label{15}
\end{equation}
is the relative concentration (particle density) of the atomic species $%
\alpha$;

\bigskip

\begin{equation}
\mathbf{j}_{i}=\pm e\sum_{\alpha}Z_{\alpha}\mathbf{\Gamma}_{\alpha} 
\label{16}
\end{equation}
is the positive ion electric current density in the boundary cells, where ''$%
+$'' is for the NORTH and EAST and ''$-$'' is for the SOUTH and WEST;

\bigskip

\begin{equation}
v_{th,e}=\sqrt{\frac{k_{B}T_{e}}{m_{e}}}   \label{17}
\end{equation}
is the electron thermal velocity;

\bigskip

\begin{equation}
J_{e}=\frac{1}{\sqrt{2\pi}}en_{e}v_{th,e}S_{1}\exp\left( -\frac{e\Phi}{%
k_{B}T_{e}}\right)   \label{18}
\end{equation}
is the electron electric current;

\bigskip

\begin{equation}
\widehat{T}_{e,i}=k_{B}T_{e,i}   \label{19}
\end{equation}
is the electron/ion temperature in the energy units and

\bigskip

\begin{equation}
\mathbf{Q}_{h\alpha D}=\kappa_{\times\alpha}\left( \mathbf{b\times\nabla }%
\widehat{T}_{\alpha}\right) S_{0}   \label{20}
\end{equation}
is the diamagnetic component of the heat current of the species $\alpha$,
where $\mathbf{b=}\frac{\mathbf{B}}{B}$.

\bigskip

\subsection{Standard Boundary Sources}

\bigskip

Analytical formulas for the standard boundary sources are presented in this
subsection. Their numerical implementation is included in the \texttt{%
b2stbc.f} source code. The B2.5 code architecture is flexible enough, so
that it is possible to change the expressions for the boundary sources
without changing the other parts of the source code, i.e., only the
numerical values of the boundary sources are used in the further
calculations.

\bigskip\ 

\subsubsection{SOUTH Boundary (Cells (ix,-1))}

\bigskip

\paragraph{Particle Sources}

\begin{equation}
S_{p\alpha}=S_{p0\alpha}+S_{p1\alpha}n_{\alpha}   \label{21}
\end{equation}

\begin{equation}
S_{p0\alpha}=C_{p0\alpha}S_{0}   \label{22}
\end{equation}

\begin{equation}
S_{p1\alpha}=C_{p1\alpha}S_{0}+C_{p2\alpha}c_{s\alpha}S_{0}+C_{p3\alpha
}c_{s\alpha}S_{1}+C_{p4\alpha}c_{st}S_{0}+C_{p5\alpha}c_{st}S_{1} 
\label{23}
\end{equation}

\bigskip

\paragraph{Momentum Sources}

\begin{equation}
S_{m\alpha}=S_{m0\alpha}+S_{m1\alpha}u_{\alpha x}+S_{m2\alpha}\rho_{\alpha
}+S_{m3\alpha}\rho_{\alpha}u_{\alpha x}   \label{24}
\end{equation}

\begin{equation}
S_{m0\alpha}=C_{m0\alpha}S_{0}   \label{25}
\end{equation}

\begin{equation}
S_{m1\alpha}=0   \label{26}
\end{equation}

\begin{equation}
S_{m2\alpha}=C_{m1\alpha}\left( -c_{s\alpha}\right) S_{1}^{\pm}+C_{m2\alpha
}\left( -c_{st}\right) S_{1}^{\pm}+C_{m5\alpha}\frac{\max\left(
0,\Gamma_{\alpha y}\right) }{n_{\alpha}}   \label{27}
\end{equation}

\begin{equation}
S_{m3\alpha}=C_{m3\alpha}S_{0}+C_{m4\alpha}S_{1}+C_{m6\alpha}\frac{%
\max\left( 0,-\Gamma_{\alpha y}\right) }{n_{\alpha}}   \label{28}
\end{equation}

\bigskip

\paragraph{All Atom/Ion Heat Source}

\begin{equation}
S_{hi}=S_{hi0}+S_{hi1}\widehat{T}_{i}+S_{hi2}n_{i}+S_{hi3}n_{i}\widehat{T}%
_{i}   \label{29}
\end{equation}

\begin{equation}
S_{hi0}=\sum_{\alpha}\left[ C_{hi0\alpha}t_{\alpha}S_{0}+C_{hi6\alpha}\max%
\left( 0,\Gamma_{\alpha y}\right) \right] +Q_{hiDy}   \label{30}
\end{equation}

\begin{equation}
S_{hi1}=\sum_{\alpha}C_{hi1\alpha}t_{\alpha}S_{0}+\frac{Q_{hiDy}}{\widehat
{T}_{i}}   \label{31}
\end{equation}

\begin{equation}
S_{hi2}=0   \label{32}
\end{equation}

\begin{equation*}
S_{hi3}=\sum_{\alpha}[C_{hi2\alpha}t_{\alpha}c_{s\alpha}S_{0}+C_{hi3\alpha
}t_{\alpha}c_{s\alpha}S_{1}+C_{hi4\alpha}t_{\alpha}c_{st}S_{0}+ 
\end{equation*}

\begin{equation}
C_{hi5\alpha}t_{\alpha}c_{st}S_{1}+C_{hi7\alpha}\frac{\max\left(
0,-\Gamma_{\alpha y}\right) }{n_{i}}]   \label{33}
\end{equation}

\bigskip

\paragraph{Electron Heat Source}

\begin{equation}
S_{he}=S_{he0}+S_{he1}\widehat{T}_{e}+S_{he2}n_{e}+S_{he3}n_{e}\widehat{T}%
_{e}   \label{34}
\end{equation}

\begin{equation}
S_{he0}=C_{he0}S_{0}+C_{he6}\max\left( 0,\Gamma_{ey}\right) +Q_{heDy} 
\label{35}
\end{equation}

\begin{equation}
S_{he1}=C_{he1}S_{0}+\frac{Q_{heDy}}{\widehat{T}_{e}}   \label{36}
\end{equation}

\begin{equation}
S_{he2}=0   \label{37}
\end{equation}

\begin{equation*}
S_{he3}=C_{he2}v_{th,e}S_{0}+C_{he3}v_{th,e}S_{1}+C_{he4}c_{st}S_{0}+ 
\end{equation*}

\begin{equation}
C_{he5}c_{st}S_{1}+C_{he7}\frac{\max\left( 0,-\Gamma_{ey}\right) }{n_{e}} 
\label{38}
\end{equation}

\bigskip

\paragraph{Electric Charge Source}

\begin{equation}
S_{ch}=S_{ch0}+S_{ch1}\Phi+S_{ch2}n_{e}+S_{ch3}n_{e}\Phi   \label{39}
\end{equation}

\begin{equation}
S_{ch0}=C_{ch7}J_{ey}-C_{ch6}J_{iy}+\max\left(
C_{ch6}J_{iy},C_{ch7}J_{ey}\right) \frac{e\Phi}{\widehat{T}_{e}}   \label{40}
\end{equation}
\bigskip

\begin{equation}
S_{ch1}=-\frac{\max\left( C_{ch6}J_{iy},C_{ch7}J_{ey}\right) e}{\widehat
{T}_{e}}   \label{41}
\end{equation}

\begin{equation}
S_{ch2}=C_{ch0}ev_{th,e}S_{0}+C_{ch2}ev_{th,e}S_{1}+C_{ch4}ec_{st}S_{1} 
\label{42}
\end{equation}

\begin{equation}
S_{ch3}=C_{ch1}\frac{e^{2}v_{th,e}S_{0}}{\widehat{T}_{e}}+C_{ch3}\frac
{e^{2}v_{th,e}S_{1}}{\widehat{T}_{e}}+C_{ch5}\frac{e^{2}c_{st}S_{1}}{%
\widehat{T}_{e}}   \label{43}
\end{equation}

\bigskip

\subsubsection{NORTH Boundary (Cells (ix,ny))}

\bigskip

\paragraph{Particle Sources}

\begin{equation}
S_{p\alpha}=S_{p0\alpha}+S_{p1\alpha}n_{\alpha}   \label{44}
\end{equation}

\begin{equation}
S_{p0\alpha}=C_{p0\alpha}S_{0}   \label{45}
\end{equation}

\begin{equation}
S_{p1\alpha}=C_{p1\alpha}S_{0}+C_{p2\alpha}c_{s\alpha}S_{0}+C_{p3\alpha
}c_{s\alpha}S_{1}+C_{p4\alpha}c_{st}S_{0}+C_{p5\alpha}c_{st}S_{1} 
\label{46}
\end{equation}

\bigskip

\paragraph{Momentum Sources}

\begin{equation}
S_{m\alpha}=S_{m0\alpha}+S_{m1\alpha}u_{\alpha x}+S_{m2\alpha}\rho_{\alpha
}+S_{m3\alpha}\rho_{\alpha}u_{\alpha x}   \label{47}
\end{equation}

\begin{equation}
S_{m0\alpha}=C_{m0\alpha}S_{0}   \label{48}
\end{equation}

\begin{equation}
S_{m1\alpha}=0   \label{49}
\end{equation}

\begin{equation}
S_{m2\alpha}=C_{m1\alpha}c_{s\alpha}S_{1}^{\pm}+C_{m2\alpha}c_{st}S_{1}^{\pm
}+C_{m5\alpha}\frac{\max\left( 0,-\Gamma_{\alpha y}\right) }{n_{\alpha}} 
\label{50}
\end{equation}

\begin{equation}
S_{m3\alpha}=C_{m3\alpha}S_{0}+C_{m4\alpha}S_{1}+C_{m6\alpha}\frac{%
\max\left( 0,\Gamma_{\alpha y}\right) }{n_{\alpha}}   \label{51}
\end{equation}

\bigskip

\paragraph{All Atom/Ion Heat Source}

\begin{equation}
S_{hi}=S_{hi0}+S_{hi1}\widehat{T}_{i}+S_{hi2}n_{i}+S_{hi3}n_{i}\widehat{T}%
_{i}   \label{52}
\end{equation}

\begin{equation}
S_{hi0}=\sum_{\alpha}\left[ C_{hi0\alpha}t_{\alpha}S_{0}+C_{hi6\alpha}\max%
\left( 0,-\Gamma_{\alpha y}\right) \right] -Q_{hiDy}   \label{53}
\end{equation}

\begin{equation}
S_{hi1}=\sum_{\alpha}C_{hi1\alpha}t_{\alpha}S_{0}-\frac{Q_{hiDy}}{\widehat
{T}_{i}}   \label{54}
\end{equation}

\begin{equation}
S_{hi2}=0   \label{55}
\end{equation}

\begin{equation*}
S_{hi3}=\sum_{\alpha}[C_{hi2\alpha}t_{\alpha}c_{s\alpha}S_{0}+C_{hi3\alpha
}t_{\alpha}c_{s\alpha}S_{1}+C_{hi4\alpha}t_{\alpha}c_{st}S_{0}+ 
\end{equation*}

\begin{equation}
C_{hi5\alpha}t_{\alpha}c_{st}S_{1}+C_{hi7\alpha}\frac{\max\left(
0,\Gamma_{\alpha y}\right) }{n_{i}}]   \label{56}
\end{equation}

\bigskip

\paragraph{Electron Heat Source}

\begin{equation}
S_{he}=S_{he0}+S_{he1}\widehat{T}_{e}+S_{he2}n_{e}+S_{he3}n_{e}\widehat{T}%
_{e}   \label{57}
\end{equation}

\begin{equation}
S_{he0}=C_{he0}S_{0}+C_{he6}\max\left( 0,-\Gamma_{ey}\right) -Q_{heDy} 
\label{58}
\end{equation}

\begin{equation}
S_{he1}=C_{he1}S_{0}-\frac{Q_{heDy}}{\widehat{T}_{e}}   \label{59}
\end{equation}

\begin{equation}
S_{he2}=0   \label{60}
\end{equation}

\begin{equation*}
S_{he3}=C_{he2}v_{th,e}S_{0}+C_{he3}v_{th,e}S_{1}+C_{he4}c_{st}S_{0}+ 
\end{equation*}

\begin{equation}
C_{he5}c_{st}S_{1}+C_{he7}\frac{\max\left( 0,\Gamma_{ey}\right) }{n_{e}} 
\label{61}
\end{equation}

\bigskip

\paragraph{Electric Charge Source}

\begin{equation}
S_{ch}=S_{ch0}+S_{ch1}\Phi+S_{ch2}n_{e}+S_{ch3}n_{e}\Phi   \label{62}
\end{equation}

\begin{equation}
S_{ch0}=C_{ch7}J_{ey}-C_{ch6}J_{iy}+\max\left(
C_{ch6}J_{iy},C_{ch7}J_{ey}\right) \frac{e\Phi}{\widehat{T}_{e}}   \label{63}
\end{equation}
\bigskip

\begin{equation}
S_{ch1}=-\frac{\max\left( C_{ch6}J_{iy},C_{ch7}J_{ey}\right) e}{\widehat
{T}_{e}}   \label{64}
\end{equation}

\begin{equation}
S_{ch2}=C_{ch0}ev_{th,e}S_{0}+C_{ch2}ev_{th,e}S_{1}+C_{ch4}ec_{st}S_{1} 
\label{65}
\end{equation}

\begin{equation}
S_{ch3}=C_{ch1}\frac{e^{2}v_{th,e}S_{0}}{\widehat{T}_{e}}+C_{ch3}\frac
{e^{2}v_{th,e}S_{1}}{\widehat{T}_{e}}+C_{ch5}\frac{e^{2}c_{st}S_{1}}{%
\widehat{T}_{e}}   \label{66}
\end{equation}

\bigskip\bigskip

\subsubsection{WEST Boundary (Cells (-1,iy), Divertor Plate)}

\bigskip

\paragraph{Particle Sources}

\begin{equation}
S_{p\alpha}=S_{p0\alpha}+S_{p1\alpha}n_{\alpha}   \label{67}
\end{equation}

\begin{equation}
S_{p0\alpha}=C_{p0\alpha}S_{0}+C_{p6\alpha}\max\left( 0,u_{\alpha
x}S_{1}^{\pm}\right)   \label{68}
\end{equation}

\begin{equation*}
S_{p1\alpha}=C_{p1\alpha}S_{0}+C_{p2\alpha}c_{sD\alpha}S_{0}+C_{p3\alpha
}c_{sD\alpha}S_{1}+C_{p4\alpha}c_{stD}S_{0}+ 
\end{equation*}

\begin{equation}
C_{p5\alpha}c_{stD}S_{1}+C_{p7\alpha}\max\left( 0,-u_{\alpha x}S_{1}^{\pm
}\right)   \label{69}
\end{equation}

\bigskip

\paragraph{Momentum Sources}

\begin{equation}
S_{m\alpha}=S_{m0\alpha}+S_{m1\alpha}u_{\alpha x}+S_{m2\alpha}\rho_{\alpha
}+S_{m3\alpha}\rho_{\alpha}u_{\alpha x}   \label{70}
\end{equation}

\begin{equation}
S_{m0\alpha}=C_{m0\alpha}S_{0}   \label{71}
\end{equation}

\begin{equation}
S_{m1\alpha}=0   \label{72}
\end{equation}

\begin{equation}
S_{m2\alpha}=C_{m1\alpha}\left( -c_{sD\alpha}\right)
S_{1}^{\pm}+C_{m2\alpha}\left( -c_{stD}\right) S_{1}^{\pm}+C_{m5\alpha}\frac
{\max\left( 0,\Gamma_{\alpha x}\right) }{n_{\alpha}}   \label{73}
\end{equation}

\begin{equation}
S_{m3\alpha}=C_{m3\alpha}S_{0}+C_{m4\alpha}S_{1}+C_{m6\alpha}\frac{%
\max\left( 0,-\Gamma_{\alpha x}\right) }{n_{\alpha}}   \label{74}
\end{equation}

\bigskip

\paragraph{All Atom/Ion Heat Source}

\begin{equation}
S_{hi}=S_{hi0}+S_{hi1}\widehat{T}_{i}+S_{hi2}n_{i}+S_{hi3}n_{i}\widehat{T}%
_{i}   \label{75}
\end{equation}

\begin{equation}
S_{hi0}=\sum_{\alpha}\left[ C_{hi0\alpha}t_{\alpha}S_{0}+C_{hi6\alpha}\max%
\left( 0,\Gamma_{\alpha x}\right) \right] +Q_{hiDx}   \label{76}
\end{equation}

\begin{equation}
S_{hi1}=\sum_{\alpha}C_{hi1\alpha}t_{\alpha}S_{0}+\frac{Q_{hiDx}}{\widehat
{T}_{i}}   \label{77}
\end{equation}

\begin{equation}
S_{hi2}=0   \label{78}
\end{equation}

\begin{equation*}
S_{hi3}=\sum_{\alpha}[C_{hi2\alpha}t_{\alpha}c_{sD\alpha}S_{0}+C_{hi3\alpha
}t_{\alpha}c_{sD\alpha}S_{1}+C_{hi4\alpha}t_{\alpha}c_{stD}S_{0}+ 
\end{equation*}

\begin{equation}
C_{hi5\alpha}t_{\alpha}c_{stD}S_{1}+C_{hi7\alpha}\frac{\max\left(
0,-\Gamma_{\alpha x}\right) }{n_{i}}]   \label{79}
\end{equation}

\bigskip

\paragraph{\protect\bigskip Electron Heat Source}

\begin{equation}
S_{he}=S_{he0}+S_{he1}\widehat{T}_{e}+S_{he2}n_{e}+S_{he3}n_{e}\widehat{T}%
_{e}   \label{80}
\end{equation}

\begin{equation}
S_{he0}=C_{he0}S_{0}+C_{he6}\max\left( 0,\Gamma_{ex}\right) +Q_{heDx} 
\label{81}
\end{equation}

\begin{equation}
S_{he1}=C_{he1}S_{0}+\frac{Q_{heDx}}{\widehat{T}_{e}}   \label{82}
\end{equation}

\begin{equation}
S_{he2}=0   \label{83}
\end{equation}

\begin{equation*}
S_{he3}=C_{he2}v_{th,e}S_{0}+C_{he3}v_{th,e}S_{1}+C_{he4}c_{stD}S_{0}+ 
\end{equation*}

\begin{equation}
C_{he5}c_{stD}S_{1}+C_{he7}\frac{\max\left( 0,-\Gamma_{ex}\right) }{n_{e}} 
\label{84}
\end{equation}

\bigskip

\paragraph{Electric Charge Source}

\begin{equation}
S_{ch}=S_{ch0}+S_{ch1}\Phi+S_{ch2}n_{e}+S_{ch3}n_{e}\Phi   \label{85}
\end{equation}

\begin{equation}
S_{ch0}=C_{ch7}J_{ex}-C_{ch6}J_{ix}+\max\left(
C_{ch6}J_{ix},C_{ch7}J_{ex}\right) \frac{e\Phi}{\widehat{T}_{e}}   \label{86}
\end{equation}
\bigskip

\begin{equation}
S_{ch1}=-\frac{\max\left( C_{ch6}J_{ix},C_{ch7}J_{ex}\right) e}{\widehat
{T}_{e}}   \label{87}
\end{equation}

\begin{equation}
S_{ch2}=C_{ch0}ev_{th,e}S_{0}+C_{ch2}ev_{th,e}S_{1}+C_{ch4}ec_{stD}S_{1} 
\label{88}
\end{equation}

\begin{equation}
S_{ch3}=C_{ch1}\frac{e^{2}v_{th,e}S_{0}}{\widehat{T}_{e}}+C_{ch3}\frac
{e^{2}v_{th,e}S_{1}}{\widehat{T}_{e}}+C_{ch5}\frac{e^{2}c_{stD}S_{1}}{%
\widehat{T}_{e}}   \label{89}
\end{equation}

\bigskip

\subsubsection{EAST Boundary (Cells (nx,iy), Divertor Plate)}

\bigskip

\paragraph{Particle Sources}

\begin{equation}
S_{p\alpha}=S_{p0\alpha}+S_{p1\alpha}n_{\alpha}   \label{90}
\end{equation}

\begin{equation}
S_{p0\alpha}=C_{p0\alpha}S_{0}+C_{p6\alpha}\max\left( 0,-u_{\alpha
x}S_{1}^{\pm}\right)   \label{91}
\end{equation}

\begin{equation*}
S_{p1\alpha}=C_{p1\alpha}S_{0}+C_{p2\alpha}c_{sD\alpha}S_{0}+C_{p3\alpha
}c_{sD\alpha}S_{1}+C_{p4\alpha}c_{stD}S_{0}+ 
\end{equation*}

\begin{equation}
C_{p5\alpha}c_{stD}S_{1}+C_{p7\alpha}\max\left( 0,u_{\alpha x}S_{1}^{\pm
}\right)   \label{92}
\end{equation}

\bigskip

\paragraph{Momentum Sources}

\begin{equation}
S_{m\alpha}=S_{m0\alpha}+S_{m1\alpha}u_{\alpha x}+S_{m2\alpha}\rho_{\alpha
}+S_{m3\alpha}\rho_{\alpha}u_{\alpha x}   \label{93}
\end{equation}

\begin{equation}
S_{m0\alpha}=C_{m0\alpha}S_{0}   \label{94}
\end{equation}

\begin{equation}
S_{m1\alpha}=0   \label{95}
\end{equation}

\begin{equation}
S_{m2\alpha}=C_{m1\alpha}c_{sD\alpha}S_{1}^{\pm}+C_{m2\alpha}c_{stD}S_{1}^{%
\pm}+C_{m5\alpha}\frac{\max\left( 0,-\Gamma_{\alpha x}\right) }{n_{\alpha }} 
\label{96}
\end{equation}

\begin{equation}
S_{m3\alpha}=C_{m3\alpha}S_{0}+C_{m4\alpha}S_{1}+C_{m6\alpha}\frac{%
\max\left( 0,\Gamma_{\alpha x}\right) }{n_{\alpha}}   \label{97}
\end{equation}

\bigskip

\paragraph{All Atom/Ion Heat Source}

\begin{equation}
S_{hi}=S_{hi0}+S_{hi1}\widehat{T}_{i}+S_{hi2}n_{i}+S_{hi3}n_{i}\widehat{T}%
_{i}   \label{98}
\end{equation}

\begin{equation}
S_{hi0}=\sum_{\alpha}\left[ C_{hi0\alpha}t_{\alpha}S_{0}+C_{hi6\alpha}\max%
\left( 0,-\Gamma_{\alpha x}\right) \right] -Q_{hiDx}   \label{99}
\end{equation}

\begin{equation}
S_{hi1}=\sum_{\alpha}C_{hi1\alpha}t_{\alpha}S_{0}-\frac{Q_{hiDx}}{\widehat
{T}_{i}}   \label{100}
\end{equation}

\begin{equation}
S_{hi2}=0   \label{101}
\end{equation}

\begin{equation*}
S_{hi3}=\sum_{\alpha}[C_{hi2\alpha}t_{\alpha}c_{sD\alpha}S_{0}+C_{hi3\alpha
}t_{\alpha}c_{sD\alpha}S_{1}+C_{hi4\alpha}t_{\alpha}c_{stD}S_{0}+ 
\end{equation*}

\begin{equation}
C_{hi5\alpha}t_{\alpha}c_{stD}S_{1}+C_{hi7\alpha}\frac{\max\left(
0,\Gamma_{\alpha x}\right) }{n_{i}}]   \label{102}
\end{equation}

\bigskip

\paragraph{\protect\bigskip Electron Heat Source}

\begin{equation}
S_{he}=S_{he0}+S_{he1}\widehat{T}_{e}+S_{he2}n_{e}+S_{he3}n_{e}\widehat{T}%
_{e}   \label{103}
\end{equation}

\begin{equation}
S_{he0}=C_{he0}S_{0}+C_{he6}\max\left( 0,-\Gamma_{ex}\right) -Q_{heDx} 
\label{104}
\end{equation}

\begin{equation}
S_{he1}=C_{he1}S_{0}-\frac{Q_{heDx}}{\widehat{T}_{e}}   \label{105}
\end{equation}

\begin{equation}
S_{he2}=0   \label{106}
\end{equation}

\begin{equation*}
S_{he3}=C_{he2}v_{th,e}S_{0}+C_{he3}v_{th,e}S_{1}+C_{he4}c_{stD}S_{0}+ 
\end{equation*}

\begin{equation}
C_{he5}c_{stD}S_{1}+C_{he7}\frac{\max\left( 0,\Gamma_{ex}\right) }{n_{e}} 
\label{107}
\end{equation}

\bigskip

\paragraph{Electric Charge Source}

\begin{equation}
S_{ch}=S_{ch0}+S_{ch1}\Phi+S_{ch2}n_{e}+S_{ch3}n_{e}\Phi   \label{108}
\end{equation}

\begin{equation}
S_{ch0}=C_{ch7}J_{ex}-C_{ch6}J_{ix}+\max\left(
C_{ch6}J_{ix},C_{ch7}J_{ex}\right) \frac{e\Phi}{\widehat{T}_{e}} 
\label{109}
\end{equation}
\bigskip

\begin{equation}
S_{ch1}=-\frac{\max\left( C_{ch6}J_{ix},C_{ch7}J_{ex}\right) e}{\widehat
{T}_{e}}   \label{110}
\end{equation}

\begin{equation}
S_{ch2}=C_{ch0}ev_{th,e}S_{0}+C_{ch2}ev_{th,e}S_{1}+C_{ch4}ec_{stD}S_{1} 
\label{111}
\end{equation}

\begin{equation}
S_{ch3}=C_{ch1}\frac{e^{2}v_{th,e}S_{0}}{\widehat{T}_{e}}+C_{ch3}\frac
{e^{2}v_{th,e}S_{1}}{\widehat{T}_{e}}+C_{ch5}\frac{e^{2}c_{stD}S_{1}}{%
\widehat{T}_{e}}   \label{112}
\end{equation}

\FRAME{fhFU}{11.3016cm}{8.4921cm}{0pt}{\Qcb{Schematic view of the
computational mesh.}}{\Qlb{fig1}}{boundcond.bmp}{\special{ language
"Scientific Word"; type "GRAPHIC"; maintain-aspect-ratio TRUE; display
"USEDEF"; valid_file "F"; width 11.3016cm; height 8.4921cm; depth 0pt;
original-width 4.2661in; original-height 3.1998in; cropleft "0"; croptop
"1"; cropright "1"; cropbottom "0"; filename 'boundcond.bmp';file-properties
"NPEU";}}\bigskip

\bigskip

\section{Standard Boundary Sources near the Material Surfaces (Divertor
Plates)}

\bigskip

\subsection{Introduction}

\bigskip

The choice of the free input parameters, which determine the standard
boundary sources, will be presented on the example of a two-component,
non-drifting ($u_{\Vert\infty}=0$) plasma (positive ions and electrons) near
the divertor plates. Ion reflection and emission from the surface are
neglected. [\ref{mriemann91}, \ref{mtakamura}] In the B2.5 code, these
boundary sources are the boundary conditions at the WEST and EAST boundary
(see Fig. \ref{fig1}). Two cases of electric current collection by the
divertor plates will be considered:

\begin{itemize}
\item  electrically floating divertor plates, which means that the net
electric current collected by the plates is zero, i.e., $I_{d}=0$, and

\item  electrically biased divertor plates, when the net electric current
collected by the plates is non-zero, i.e., $I_{d}\neq0$.
\end{itemize}

\bigskip

In the second case it is assumed that the net current does not cause some
macroscopic current-driven plasma instabilities and that the spatial
variation of the electric potential near the wall is qualitatively similar
as in the floating case.

\bigskip

\subsection{Standard Particle Sources}

\bigskip

The poloidal component of the ion flux at the magnetic presheath edge ($MPSE$%
) is given by:

\begin{equation}
\Gamma_{xMPSE}=n_{MPSE}\left\langle v_{x}\right\rangle _{MPSE}\text{,} 
\label{113}
\end{equation}
where $n_{MPSE}$ is the ion particle density and $\left\langle
v_{x}\right\rangle _{MPSE}$ is the poloidal component of the ion drift
velocity, both calculated at the magnetic presheath edge. According to the
fluid model [\ref{mstangeby95}], equation (\ref{113}) can be rewritten as:

\begin{equation}
\Gamma_{xMPSE}\circeq n_{MPSE}u_{\Vert MPSE}\sin\theta\text{,}   \label{113a}
\end{equation}
\bigskip where $\sin\theta=\frac{B_{x}}{B}$ and

\begin{equation}
u_{\Vert MPSE}=c_{sMPSE}\pm\mathcal{O}\left( \frac{u_{DMPSE}}{\tan\theta }%
\right) \text{,}   \label{114}
\end{equation}
where $c_{sMPSE}=\sqrt{\frac{k_{B}(T_{eMPSE}+T_{iMPSE})}{m_{i}}}$ is the
local ion sound velocity and $u_{DMPSE}=\left( \frac{E_{y}}{B}\right) _{MPSE}
$ is the local drift velocity due to the poloidal $\mathbf{E\times B}$ drift
in the magnetic presheath (in the B2.5 geometry this drift is caused by the
radial electric field $E_{y}$). Equations (\ref{11}) and (\ref{12}) for
conditionally corrected ion drift velocities $c_{sD\alpha}$ and $c_{stD}$
(the numerical factor $0.1$ is due to a specific value of $\sin\theta$, for
which the value of the ion drift velocity should be corrected) are based on
the approximation:

\begin{equation}
u_{\Vert MPSE}\circeq c_{sMPSE}\pm\frac{u_{DMPSE}}{\tan\theta}\text{.} 
\label{115}
\end{equation}
Better approximation and more complicated analytical formula for $u_{\Vert
MPSE}$ can be found in the article of P. C. Stangeby and A. V. Chankin [\ref
{mstangeby95}] (the coordinate system in their calculation is different, $x$
is the radial component and $y$ is the poloidal component). According to
their model, the linear approximation of the expression for $u_{\Vert MPSE}$
is:

\begin{equation}
u_{\Vert MPSE}\circeq c_{sMPSE}-\frac{l_{sh}}{2\left| \Psi_{W}\right|
\tan\theta}\left( \frac{1}{l_{v}}-\frac{1}{l_{n}}\right) u_{DMPSE}\text{,} 
\label{stanchan}
\end{equation}
where $\Psi_{W}$ is normalized sheath electric potential fall (typically $%
\left| \Psi_{W}\right| \approx3$ for hydrogenic plasmas) and $l_{sh}$, $l_{v}
$ and $l_{n}$ are characteristic decay lengths for the sheath electric
potential, ion drift velocity and particle density respectively. [\ref
{mstangeby95}] Assumption $l_{v}=l_{n}$ in equation (\ref{stanchan}) gives
zero contribution of the electric drift term, which is probably physically
unrealistic, because it means that even very large poloidal $\mathbf{E\times
B}$ drifts, so that local $c_{s}\approx\frac{E_{y}}{B}$, do not contribute
in the linear order to the ion drift velocity at the magnetic presheath
edge. Implementation of their non-linear expression for $u_{\Vert MPSE}$
into the B2.5 code demands significant changes of the standard source code.

\bigskip

The alternative approach is to specify the ion flux at the electrostatic
sheath edge:

\begin{equation}
\Gamma_{xSE}=n_{SE}\left\langle v_{x}\right\rangle _{SE}\text{,} 
\label{sheathFlux}
\end{equation}
which is equal to the ion flux at the magnetic presheath edge $\Gamma_{xMPSE}
$. In this case, it can be assumed that:

\begin{equation}
\left\langle v_{x}\right\rangle _{SE}=c_{sSE}\sin\theta\text{.} 
\label{sheathVel}
\end{equation}
However, the ion density at the electrostatic sheath edge $n_{SE}$ is
different from the ion density at the magnetic presheath edge $n_{MPSE}$ and
it can be some very complicated non-linear function of other plasma
parameters. In order to apply the expressions for the particle, momentum and
energy fluxes at the electrostatic sheath edge, it is necessary to develop a
kinetic model of the collisionless plasma region near the solid surface,
because the fluid models cannot describe correctly this region. In
principle, the validity of the B2.5 equations is limited only to the plasma
regions where the ion velocity distributions do not differ significantly
from Maxwellian or drifting-Maxwellian distributions, which is not the case
also in the plasma presheath. The correct two-scale or multi-scale analysis
of the boundary plasma near the divertor plate should include a separate
description of the collisional presheath, which is located between the bulk
plasma and collisionless magnetic presheath. [\ref{mriemann91}] Another
possibility is to develop a model of a complete plasma presheath, which
includes the collisional presheath and collisionless magnetic presheath. In
any case, the solutions of the equations for these different plasma regions
should be smoothly connected and implemented into the boundary conditions of
the B2.5 code.

\bigskip

In the standard version of the boundary sources, the particle source can be
specified by choosing $C_{p3}=1$ and all other $C_{pj}=0$. The numerical
factor $0.1$ in the expressions for $c_{sD\alpha}$ and $c_{stD}$ (see
equations (\ref{11}) and (\ref{12})) can be replaced by $\sin\theta$, so
that they can be more general. Another possibility is to redefine $%
c_{sD\alpha}$ given by equation (\ref{11}), so that the linearized
expression for the particle flux of species $\alpha$ will be:

\begin{equation}
\Gamma_{\alpha xMPSE}\circeq n_{\alpha MPSE}\left( c_{s\alpha MPSE}\sin
\theta\pm a_{D}u_{DMPSE}\cos\theta\right) \text{,}   \label{116}
\end{equation}
where the additional numerical factor $a_{D}$ is calculated with some
theoretical model of the magnetic presheath (e.g., [\ref{mstangeby95}]). In
the absence of any reliable theoretical model it can be assumed that $a_{D}=0
$ or $a_{D}=1$, which corresponds to the electric drift-free case and to the
approximation (\ref{115}) respectively. If these small changes of the source
code are applied, the appropriate choice of the free input parameters for
the particle sources will be again $C_{p3\alpha}=1$ and all other $%
C_{pj\alpha}=0$.

\bigskip

However, many experiments (see [\ref{mmatt90}], [\ref{mstanpitch}], [\ref
{mmatt94}], [\ref{mpitch97}])$\ $and some numerical simulations [\ref{mpic89}%
] have shown that when the angle $\theta$ is very small, i.e., especially
for $\theta\sim\sqrt{\frac{m_{e}}{m_{i}}}$, which is in practical cases
usually for $\theta\lesssim5^{\circ}$, the expression for the ion flux (\ref
{113a}) is not correct. In such cases, the measured charged particle fluxes
to the surface are anomalously high compared with the theoretical prediction
given by equation (\ref{113a}). They can be approximately described with a
turbulent diffusion model:

\begin{equation}
\Gamma_{xMPSE}\circeq-D_{turb}\partial_{x}n_{MPSE}\approx\frac{D_{turb}}{%
\lambda_{MPSx}}n_{MPSE}\text{,}   \label{117}
\end{equation}
where $D_{turb}$ is the local turbulent diffusion coefficient and $%
\lambda_{MPSx}$ is the characteristic poloidal dimension of the magnetic
presheath. The quotient $\frac{D_{turb}}{\lambda_{MPSx}}$ has a dimension of
a velocity and the approximation of the right-hand side of equation (\ref
{117}) enables its implementation into the B2.5 code by using the
expressions for the boundary sources. In principle, only the theoretical
model of turbulent plasma transport for this specific case can give a
correct result for the coefficient $S_{p1}$(see equations (\ref{69}) and (%
\ref{92})). In the absence of such a theoretical model it is necessary to
estimate the quotient $\frac{D_{turb}}{\lambda_{MPSx}}$ properly, so that
the numerical values of the free input parameters for the particle sources
can be chosen.

\bigskip

According to the fluid models of the magnetized boundary plasma, the
characteristic length scale of the magnetic presheath is $\rho_{s}\cos\theta$%
, where $\rho_{s}=\frac{c_{s}}{\omega_{i}}$ is the Larmor radius for the
ions moving with the ion sound velocity. Since in our case $\theta\sim\sqrt {%
\frac{m_{e}}{m_{i}}}\ll1$, $\lambda_{MPSx}$ is of the order of magnitude of $%
\rho_{sMPSE}$, i.e., $\lambda_{MPSx}\sim\rho_{sMPSE}$. The turbulent
diffusion coefficient $D_{turb}$ is approximately given by:

\begin{equation}
D_{turb}\approx\sqrt{\left\langle \widetilde{u}_{D}^{2}\right\rangle }%
\lambda_{c}\text{,}   \label{118}
\end{equation}
where $\widetilde{u}_{D}=\frac{\widetilde{E}_{y}}{B}$ is the drift velocity
due to the fluctuating (turbulent) electric field $\widetilde{E}_{y}$ and $%
\lambda_{c}$ $\approx\left\langle k_{y}\right\rangle ^{-1}$ is the
characteristic coherence length of the turbulent spectrum. Turbulent
fluctuations of the magnetic field are much smaller and can be neglected in
the study of the particle transport. (However, they can be important for the
heat transport. [\ref{mconnor}]) Assuming $\sqrt{\left\langle \widetilde
{E}_{y}^{2}\right\rangle }\approx\sqrt{\left\langle \widetilde{\Phi}%
^{2}\right\rangle }\left\langle k_{y}\right\rangle \equiv\delta\Phi
\left\langle k_{y}\right\rangle $, equation (\ref{118}) can be rewritten as:

\begin{equation}
D_{turb}\approx\frac{\delta\Phi}{B}\text{.}   \label{119}
\end{equation}
Turbulent fluctuations of the electric potential $\delta\Phi$ are
approximately proportional to the turbulent fluctuations of the plasma
density, i.e.,
\begin{equation}
\frac{\delta n}{\left\langle n\right\rangle }\approx\frac{e\delta\Phi}{%
k_{B}T_{e}}\text{,}   \label{120}
\end{equation}
where so-called Boltzmann's fluctuations are assumed and possible
decorrelations between the potential and density fluctuations are neglected.
In the electrostatic turbulence, the fluctuations of the electron
temperature are usually much smaller than the density and potential
fluctuations, so that they can be also neglected. Equations (\ref{119}) and (%
\ref{120}) can be combined to give:

\begin{equation}
D_{turb}\approx\frac{\delta n}{\left\langle n\right\rangle }\frac{k_{B}T_{e}%
}{eB}\text{.}   \label{121}
\end{equation}
This simple calculation gives the characteristic Bohm $\frac{1}{B}$-scaling
of the anomalous diffusion coefficient. In the original Bohm diffusion
coefficient, there is an empirical numerical factor $\frac{1}{16}$ instead
of $\frac{\delta n}{\left\langle n\right\rangle }$. The physical picture of
the result (\ref{121}) is clear: increased turbulence level results in
increased turbulent transport of the charged particles. Assuming $%
\lambda_{MPSx}\sim \rho_{sMPSE}$, the quotient $\frac{D_{turb}}{%
\lambda_{MPSx}}$ is approximately:

\begin{equation}
\frac{D_{turb}}{\lambda_{MPSx}}\approx\left( \frac{\delta n}{\left\langle
n\right\rangle }\right) _{MPSE}c_{sMPSE}\text{.}   \label{122}
\end{equation}
This result qualitatively explains the anomalously high particle fluxes at
small grazing angles $\theta$. Despite the fact that collisions in the
divertor plasma can significantly reduce the turbulent fluctuation level, it
is still valid that $\frac{\delta n}{\left\langle n\right\rangle }\gg
\sin\theta$ in such cases. If the information about the density fluctuation
level $\frac{\delta n}{\left\langle n\right\rangle }$ is available, e.g.,
from the experiment, this simple result can be used in the expression for
the particle source by choosing $C_{p2}=\left( \frac{\delta n}{\left\langle
n\right\rangle }\right) _{MPSE}$ and all other $C_{pj}=0$. The other
possibility is to calculate $\frac{D_{turb}}{\lambda_{MPSx}}$ somehow and
choose $C_{p1}=\frac{D_{turb}}{\lambda_{MPSx}}$ and all other $C_{pj}=0$. It
is also important to mention that the turbulent transport near the material
wall can be different from the turbulent transport in the bulk plasma, which
means that the corresponding macroscopic turbulent transport coefficients
can be different.

\bigskip

Equations (\ref{113a}) and (\ref{117}) may be considered as limiting cases,
because in the general case or for some values of the grazing angle $\theta$
the classical and turbulent contributions to the plasma transport can be
comparable by magnitude. In the bulk plasma transport there are also
contributions due to the classical and turbulent transport. Thus it is
reasonable to assume the following expression for the ion flux at the
magnetic presheath edge:

\begin{equation}
\Gamma_{xMPSE}\circeq n_{MPSE}\left( c_{sMPSE}\sin\theta\pm
a_{D}u_{DMPSE}\cos\theta\right) -D_{turb}\partial_{x}n_{MPSE}\text{.} 
\label{totalGamma}
\end{equation}
Equation (\ref{totalGamma}) cannot be directly implemented in the standard
expressions for the boundary sources, because they do not include
derivatives of the plasma parameters, such as $\partial_{x}n_{\alpha}$, $%
\partial _{x}\widehat{T}_{i}$, $\partial_{x}\widehat{T}_{e}$ etc. There are
two practical solutions for this problem. The first solution is to add in
the source code new terms with the derivatives of the plasma parameters in
the expressions for the boundary sources, so that equation (\ref{totalGamma}%
) can be implemented directly. The second solution is to estimate the
turbulent component of the ion flux, as it was shown in the case of a very
small grazing angle $\theta$, so that the appropriate choices of the free
input parameters will be $C_{p1}=\frac{D_{turb}}{\lambda_{MPSx}}$, $C_{p3}=1$
and all other $C_{pj}=0$, or $C_{p2}=\left( \frac{\delta n}{\left\langle
n\right\rangle }\right) _{MPSE}$, $C_{p3}=1$ and all other $C_{pj}=0$.

\bigskip

\subsection{Standard Momentum Sources}

\bigskip

In order to calculate the boundary sources for the momentum conservation
equation, for $d\mathbf{A}=dA\mathbf{e}_{x}$ equation (\ref{2}) can be
rewritten as:

\begin{equation}
\mathbf{S}_{m}=-\int_{B.S.}\underline{\mathbf{M}}\cdot\mathbf{e}%
_{x}dA\equiv-\int_{B.S.}\mathbf{P}_{m}dA\text{,}   \label{123}
\end{equation}
where the momentum flux $\mathbf{P}_{m}$ is given by [\ref{manders90}]:

\begin{equation}
\mathbf{P}_{m}=m_{i}n\left[ \left\langle v_{x}^{2}\right\rangle ,\text{ }%
\left\langle v_{y}v_{x}\right\rangle ,\text{ }\left\langle
v_{z}v_{x}\right\rangle \right] ^{T}\text{.}   \label{124}
\end{equation}
All quantities in equation (\ref{124}) should be calculated at the magnetic
presheath edge ($MPSE$). The velocity $\mathbf{v}$ may be decomposed into an
average part $\mathbf{u}$ and a random part $\mathbf{w}$, associated with a
``thermal'' motion, i.e.,

\begin{equation}
\mathbf{v=u+w}\text{,}   \label{125}
\end{equation}
where $\left\langle \mathbf{w}\right\rangle =0$. Then the components of the
momentum flux $\mathbf{P}_{m}$ can be written as:

\begin{equation}
P_{mx}=m_{i}n\left( u_{x}^{2}+\left\langle w_{x}^{2}\right\rangle \right) 
\text{,}   \label{126}
\end{equation}

\begin{equation}
P_{my}=m_{i}n\left( u_{y}u_{x}+\left\langle w_{y}w_{x}\right\rangle \right) 
\text{,}   \label{127}
\end{equation}

\begin{equation}
P_{mz}=m_{i}n\left( u_{z}u_{x}+\left\langle w_{z}w_{x}\right\rangle \right) 
\text{.}   \label{128}
\end{equation}

\bigskip

In further derivations some assumptions about the ion velocity distribution
at the magnetic presheath edge should be made. The assumption that the ion
velocity distribution is drifting-Maxwellian is acceptable from the point of
view of the bulk plasma transport, but from the point of view of the
two-scale analysis of the plasma-sheath problem such an ion velocity
distribution cannot satisfy the generalized Bohm criterion [\ref{mriemann91}%
]. In the absence of any better approximation and assuming that $%
u_{z}\approx0$ and $u_{y}\approx 0$, the only non-zero component of $\mathbf{%
P}_{m}$ is $P_{mx}$, because for the Maxwellian distribution the anisotropic
shear (viscosity) terms $m_{i}n\left\langle w_{y}w_{x}\right\rangle $ and $%
m_{i}n\left\langle w_{z}w_{x}\right\rangle $ are zero. This assumption is
maybe acceptable for the component $P_{mz}$, but it may be very questionable
for the component $P_{my}$, because, according to the study of A. V. Chankin
and P. C. Stangeby [\ref{mchankin94}], the radial component of the
diamagnetic drift velocity:

\begin{equation}
\mathbf{u}_{Dm}=\frac{\mathbf{B\times\nabla} p}{enB^{2}}\text{,} 
\label{129}
\end{equation}
which is in our case:

\begin{equation}
u_{Dmy}=\frac{B_{z}\partial_{x}p-B_{x}\partial_{z}p}{enB^{2}}\text{,} 
\label{130}
\end{equation}
and the corresponding particle flux $\Gamma_{Dmy}=nu_{Dmy}$ parallel to the
material surface can be significant. It also means that the $v_{y}$ ion
velocity distribution at the magnetic presheath edge may be non-Maxwellian,
which gives a non-zero contribution of the viscosity term $%
m_{i}n\left\langle w_{y}w_{x}\right\rangle $. In addition to that, it should
be mentioned that the radial viscosity in the bulk plasma is predominantly
turbulent, so it is very likely that it is quite significant also at the
magnetic presheath edge. However, in the standard version of the B2.5 code
these contributions to the momentum fluxes at the material surfaces are not
taken into account and even $P_{mx}$ is calculated very approximately. Thus
neglecting $P_{my}$ and $P_{mz}$, the momentum flux at the magnetic
presheath edge can be written as:

\begin{equation}
P_{mxMPSE}\circeq\left( \rho_{MPSE}u_{\Vert MPSE}^{2}+p_{iMPSE}\right)
\sin\theta-m_{i}c_{sMPSE}D_{turb}\partial_{x}n_{MPSE}\text{,}   \label{131}
\end{equation}
where $\rho_{MPSE}=m_{i}n_{MPSE}$ is the ion mass density and $%
p_{iMPSE}=n_{MPSE}k_{B}$ $T_{iMPSE}$ is the ion pressure, both calculated at
the magnetic presheath edge. The last term on the right-hand side of
equation (\ref{131}) is a phenomenological turbulent momentum flux. In the
case of a very small grazing angle $\theta$, this term is much larger than
the classical term. Comparison of equation (\ref{131}) with equations (\ref
{70}) and (\ref{93}) shows that equation (\ref{131}) cannot be implemented
directly in the standard expressions for the momentum sources. It includes
more complicated non-linear dependence on the plasma parameters than
equations for the momentum sources (\ref{70}) and (\ref{93}), so that they
cannot accurately describe the boundary values of the momentum fluxes. It is
necessary to make some additional approximations in order to implement
equation (\ref{131}) in the standard expressions for the momentum sources.
Assuming that $T_{iMPSE}\approx T_{eMPSE}$, $u_{\Vert MPSE}\approx c_{sMPSE}$
and using the same approximation for the turbulent transport as in the upper
analysis of the particle flux, the momentum flux at the magnetic presheath
edge can be written approximately as:

\begin{equation}
P_{mxMPSE}\approx\frac{3}{2}\rho_{MPSE}u_{\Vert MPSE}^{2}\sin\theta+\left( 
\frac{\delta n}{\left\langle n\right\rangle }\right) _{MPSE}\rho
_{MPSE}c_{sMPSE}^{2}\text{.}   \label{132}
\end{equation}

\bigskip

There are several possible choices of the free input parameters $C_{mj}$ and
all of them give very approximate values of the boundary sources. The first
possibility is to choose $C_{m0}=P_{mxMPSE}$, where $P_{mxMPSE}$ is
calculated by using the initial or some other (e.g. experimental) specific
values of the ion density, electron and ion temperature, and all other $%
C_{mj}=0$. This possibility will be probably the best one in the case of a
very small grazing angle $\theta$, i.e., for $\theta\lesssim5^{\circ}$, if
the value of the momentum flux is provided from some turbulent transport
model (see, e.g., [\ref{mconnor}] and references therein). The second
possibility is to specify the free input parameters in the expression for $%
S_{m2}$ (see equations (\ref{73}) and (\ref{96})). In this case, assuming
that $u_{\Vert MPSE}\approx c_{sMPSE}$ the possible choice is $C_{m1}=\left[ 
\frac{3}{2}+\frac{\left( \frac{\delta n}{\left\langle n\right\rangle }%
\right) _{MPSE}}{\sin\theta }\right] c_{sMPSE}$ and all other $C_{mj}=0$.
The third possibility is to specify the free input parameters in the
expression for $S_{m3}$ (see equations (\ref{74}) and (\ref{97})). The fluid
velocities $u_{\alpha x}$ in equations (\ref{70}) and (\ref{93}) are
calculated with the bulk transport equations, which means that, in order to
preserve the validity of the fluid model and to a certain extent assure the
stability of the ion fluid flow, they should not exceed the thermal
velocities, i.e., $u_{\alpha x}\leq v_{th,\alpha}\sin\theta$. Then, assuming
that $u_{\Vert MPSE}^{2}\approx c_{sMPSE}^{2}\approx2v_{th,i}^{2}$, the
appropriate choice of the free input parameters can be $C_{m3}=\left[ \frac{3%
}{\sqrt{2}}+\left( \frac{\delta n}{\left\langle n\right\rangle }\right)
_{MPSE}\frac{\sqrt{2}}{\sin\theta }\right] c_{sMPSE}$ and all other $C_{mj}=0
$. In all these cases the initial or some other specific values of the
plasma parameters should be used in the calculations of the free input
parameters $C_{mj}$.

\bigskip

\subsection{Standard Ion and Electron Heat Sources}

\bigskip

The ion heat flux at the magnetic presheath edge, $Q_{ix}=\frac{1}{2}%
m_{i}n\left\langle v^{2}v_{x}\right\rangle $, can be written as (index $MPSE$
is dropped) [\ref{manders90}]:

\begin{equation}
Q_{ix}=\frac{1}{2}m_{i}nu^{2}u_{x}+\frac{5}{2}p_{i}u_{x}+\Pi_{xx}u_{x}+%
\Pi_{xy}u_{y}+\Pi_{xz}u_{z}+q_{ix}\text{,}   \label{133}
\end{equation}
where

\begin{equation}
p_{i}=\frac{1}{3}m_{i}n\left\langle w^{2}\right\rangle =nk_{B}T_{i} 
\label{134}
\end{equation}
is the scalar pressure,

\begin{equation}
\Pi_{jk}=m_{i}n\left\langle w_{j}w_{k}-\frac{1}{3}w^{2}\delta_{jk}\right%
\rangle   \label{135}
\end{equation}
is the anisotropic shear (viscosity) tensor and

\begin{equation}
q_{ix}=\frac{1}{2}m_{i}n\left\langle w^{2}w_{x}\right\rangle   \label{136}
\end{equation}
is the heat conduction flux. All these quantities should be calculated at
the magnetic presheath edge. It is obvious again, as in the analysis of the
momentum flux, that the contribution of the non-zero radial drift velocity $%
u_{y}$, i.e., the term $\Pi_{xy}u_{y}$, in equation (\ref{133}) may be
important. Neglecting terms with $u_{y}$ and $u_{z}$ in equation (\ref{133})
and assuming again that the ion velocity distribution at the magnetic
presheath edge is drifting-Maxwellian, the ion heat flux at the magnetic
presheath edge can be written approximately as:

\begin{equation*}
Q_{ixMPSE}\circeq{\Huge (}\frac{1}{2}m_{i}n_{MPSE}u_{\Vert MPSE}^{3}+\frac
{5}{2}n_{MPSE}k_{B}T_{iMPSE}u_{\Vert MPSE}
\end{equation*}

\begin{equation}
+2k_{B}T_{iMPSE}\Gamma_{th,iMPSE}{\Huge )}\sin\theta-\kappa_{turb}^{i}%
\partial_{x}\widehat{T}_{iMPSE}\text{,}   \label{137}
\end{equation}
\bigskip where $\Gamma_{th,iMPSE}=n_{MPSE}\sqrt{\frac{k_{B}T_{iMPSE}}{2\pi
m_{i}}}$ is the ion thermal flux at the magnetic presheath edge. The last
term on the right-hand side of equation (\ref{137}) is a phenomenological
turbulent ion heat conduction flux, which is much larger than the classical
ion heat flux, when the grazing angle $\theta$ is very small. The ion heat
flux at the material wall, $Q_{ixW}$, is larger because of the electric
potential drop between the magnetic presheath edge and the wall, so that:

\begin{equation}
Q_{ixW}=Q_{ixMPSE}+e\left( \Phi_{MPSE}-\Phi_{W}\right) \Gamma_{xMPSE}\text{.}
\label{138}
\end{equation}
On the other hand, this potential drop reduces the electron heat flux at the
wall, $Q_{exW}$, so that:

\begin{equation}
Q_{exW}=Q_{exMPSE}-e\left( \Phi_{MPSE}-\Phi_{W}\right) \Gamma_{xMPSE}\text{.}
\label{139}
\end{equation}
The electron heat flux at the wall, $Q_{exW}$, is given by:

\begin{equation*}
Q_{exW}\circeq\sqrt{\frac{2}{\pi}}n_{eMPSE}v_{th,eMPSE}k_{B}T_{eMPSE}\exp%
\left[ \frac{e\left( \Phi_{W}-\Phi_{MPSE}\right) }{k_{B}T_{eMPSE}}\right]
\sin\theta 
\end{equation*}

\begin{equation}
-\kappa_{turb}^{e}\partial_{x}\widehat{T}_{eMPSE}\text{,}   \label{140}
\end{equation}
where $n_{eMPSE}=n_{MPSE}$ in a case of a two-component plasma. The second
term on the right-hand side of equation (\ref{140}) is a phenomenological
turbulent electron heat conduction flux, which is much larger than the
classical electron heat flux, when the grazing angle $\theta$ is very small.
Since the boundary sources include the collisionless region with the
electric potential drop, its effect on the charged particle heat fluxes must
be taken into account. Summing of equations (\ref{138}) and (\ref{139})
gives:

\begin{equation}
Q_{ixW}+Q_{exW}=Q_{ixMPSE}+Q_{exMPSE}\text{,}   \label{141}
\end{equation}
which is just an energy conservation law. In the expressions for the
boundary heat sources it is possible to use charged particle heat fluxes
calculated at the material wall or at the magnetic presheath edge, but in
any case the energy conservation law (\ref{141}) must be satisfied. In this
example, the expressions for the boundary heat sources will be calculated
with the charged particle heat fluxes calculated at the material wall.

\bigskip

Expressions for the ion and electron heat flux (\ref{137})-(\ref{140})
include more complicated non-linear dependence on the plasma parameters than
the standard expressions for the boundary heat sources presented in the
previous section (see equations (\ref{75}), (\ref{80}), (\ref{98}) and (\ref
{103})). Thus additional approximations in equations (\ref{137})-(\ref{140})
are necessary in order to apply them in the standard expressions for the
boundary heat sources.

\bigskip

In the absence of the turbulent transport model it is very difficult to make
good approximations for the turbulent ion and electron heat flux. In both
cases it can be assumed that (indices $i$, $e$ and $MPSE$ are dropped):

\begin{equation}
Q_{x(turb)}\approx\frac{\kappa_{turb}}{\lambda_{Tx}}k_{B}T\approx2k_{B}T%
\Gamma_{xturb}\approx2k_{B}T\frac{D_{turb}}{\lambda_{nx}}n\text{,} 
\label{142}
\end{equation}
where $\lambda_{Tx}$ and $\lambda_{nx}$ are characteristic decay lengths for
the temperature and density respectively. Assuming $\lambda_{Tx}\approx 
\frac{\lambda_{nx}}{2}\approx\frac{\lambda_{MPSx}}{2}$, which gives $%
\kappa_{turb}\approx nD_{turb}$, and with the upper estimate for $\frac{%
D_{turb}}{\lambda_{MPSx}}$ (see equation (\ref{122})), a relatively simple
expression for the turbulent heat flux is obtained:

\begin{equation}
Q_{i,exMPSE(turb)}\approx2\left( \frac{\delta n}{\left\langle n\right\rangle 
}\right) _{MPSE}c_{sMPSE}n_{MPSE}k_{B}T_{i,eMPSE}\text{.}   \label{turbQ}
\end{equation}
Approximate expression for the turbulent heat flux (\ref{turbQ}) can be
applied in the standard expression for the boundary heat source. Assuming $%
T_{iMPSE}\approx T_{eMPSE}$ and $u_{\Vert MPSE}\approx c_{sMPSE}\approx 
\sqrt{2}v_{th,iMPSE}$ in the classical ion heat flux, the approximate
expression for the ion heat flux at the wall is:

\begin{equation*}
Q_{ixW}\approx\left[ \frac{7}{2}+\frac{1}{\sqrt{\pi}}+\frac{e\left(
\Phi_{MPSE}-\Phi_{W}\right) }{k_{B}T_{eMPSE}}\right] n_{MPSE}u_{\Vert
MPSE}k_{B}T_{iMPSE}\sin\theta 
\end{equation*}

\begin{equation}
+\left[ 2+\frac{e\left( \Phi_{MPSE}-\Phi_{W}\right) }{k_{B}T_{eMPSE}}\right]
\left( \frac{\delta n}{\left\langle n\right\rangle }\right)
_{MPSE}n_{MPSE}c_{sMPSE}k_{B}T_{iMPSE}\text{.}   \label{approxQi}
\end{equation}
The approximate expression for the electron heat flux at the wall is:

\begin{equation*}
Q_{exW}\approx\sqrt{\frac{2}{\pi}}n_{eMPSE}v_{th,eMPSE}k_{B}T_{eMPSE}\exp%
\left[ \frac{e\left( \Phi_{W}-\Phi_{MPSE}\right) }{k_{B}T_{eMPSE}}\right]
\sin\theta 
\end{equation*}

\begin{equation}
+2\left( \frac{\delta n}{\left\langle n\right\rangle }\right)
_{MPSE}n_{MPSE}c_{sMPSE}k_{B}T_{eMPSE}\text{.}   \label{approxQe}
\end{equation}

\bigskip

In order to specify the numerical values of the free input parameters $%
C_{hij}$ and $C_{hej}$, it is also necessary to obtain the numerical value
for the normalized potential fall $\frac{e\left( \Phi_{W}-\Phi_{MPSE}\right) 
}{k_{B}T_{eMPSE}}$ in equations (\ref{approxQi}) and (\ref{approxQe}). It
can be obtained directly from the experiment, e.g., from measurements of the
spatial profiles of the electric potential with emissive probes and electron
temperature with Langmuir (collecting) probes, or calculated from the
measured electric current to the divertor plate, $I_{d}=\pm j_{d}A_{d}$. The
total electric current density at the divertor plate, $j_{d}$, can be
written as:

\begin{equation*}
j_{d}=j_{iMPSEx}+j_{eMPSEx}\circeq en_{MPSE}u_{\Vert MPSE}\sin\theta- 
\end{equation*}

\begin{equation}
\frac{\left( 1-\gamma_{se}\right) }{\sqrt{2\pi}}en_{eMPSE}v_{th,eMPSE}\exp%
\left[ \frac{e\left( \Phi_{W}-\Phi_{MPSE}\right) }{k_{B}T_{eMPSE}}\right]
\sin\theta\text{,}   \label{jdTotal}
\end{equation}
where $\gamma_{se}$ is the secondary electron emission coefficient and it is
assumed that the turbulent transport is ambipolar, i.e., $%
D_{turb}^{i}=D_{turb}^{e}$, so that its net contribution to the total
electric current is zero. The secondary electron emission is not taken into
account in the expression for the electron current (\ref{18}), although it
can be important in some practical cases. In such cases it is necessary to
include the correction due to the secondary electron emission in the
expression for the electric charge source, which is presented in the text
below. If the secondary electron emission varies along the material surface,
so that the effective secondary electron emission coefficient is not
constant, this effect must be explicitly taken into account in the
expression for the electron current (\ref{18}). Equation (\ref{jdTotal}) is
applicable only if the grazing angle $\theta$ is not very small, i.e., if $%
\theta\gtrsim5^{\circ}$. In the case of a very small grazing angle $\theta$
it may be assumed that $\frac{e\left( \Phi_{W}-\Phi_{MPSE}\right) }{%
k_{B}T_{eMPSE}}\approx0$ in equations (\ref{approxQi}) and (\ref{approxQe})
or the experimental value for that specific case should be used. Considering 
$n_{eMPSE}=n_{MPSE}$, equation (\ref{jdTotal}) can be rewritten as:

\begin{equation*}
\frac{e\left( \Phi_{W}-\Phi_{MPSE}\right) }{k_{B}T_{eMPSE}}= 
\end{equation*}

\begin{equation}
\ln\left[ \frac{\sqrt{2\pi}u_{\Vert MPSE}}{\left( 1-\gamma_{se}\right)
v_{th,eMPSE}}-\frac{j_{d}}{\frac{\left( 1-\gamma_{se}\right) }{\sqrt{2\pi}}%
en_{MPSE}v_{th,eMPSE}\sin\theta}\right] \text{,}   \label{deltaPhi}
\end{equation}
where some specific values of the plasma parameters, e.g., the initial
values in the simulation ($n_{i,e}^{\left( 0\right) }$, $\widehat{T}%
_{i,e}^{\left( 0\right) }$, $u_{D}^{\left( 0\right) }$) or experimental
values, should be used in the right-hand side of equation (\ref{deltaPhi}).
For example, it can be assumed that $u_{\Vert MPSE}\approx c_{sMPSE}$ and
for the plasma density at the magnetic presheath edge $n_{MPSE}\approx
n_{i,e}^{\left( 0\right) }\exp\left( \frac{e\left(
\Phi_{MPSE}-\Phi_{\infty}^{\left( 0\right) }\right) }{k_{B}T_{e}^{\left(
0\right) }}\right) \approx n_{i,e}^{\left( 0\right) }\exp\left( -\frac{1}{2}%
\right) $, because the potential fall in the presheath, which is
approximately $-\frac{k_{B}T_{e}^{\left( 0\right) }}{2e}$, accelerates the
ions from their thermal velocities up to the ion sound velocity. In a case
of electrically floating divertor plates, $j_{d}=0$, so that:

\begin{equation}
\frac{e\left( \Phi_{W}^{f}-\Phi_{MPSE}^{f}\right) }{k_{B}T_{eMPSE}}=\ln\left[
\frac{\sqrt{2\pi}u_{\Vert MPSE}}{\left( 1-\gamma_{se}\right) v_{th,eMPSE}}%
\right] \text{.}   \label{floatPhi}
\end{equation}
Assuming $u_{\Vert MPSE}\approx c_{sMPSE}$, equation (\ref{floatPhi}) can be
rewritten as:

\begin{equation}
\frac{e\left( \Phi_{W}^{f}-\Phi_{MPSE}^{f}\right) }{k_{B}T_{eMPSE}}\circeq%
\frac{1}{2}\ln\left[ \frac{2\pi m_{e}\left( 1+\frac{T_{iMPSE}}{T_{eMPSE}}%
\right) }{m_{i}\left( 1-\gamma_{se}\right) ^{2}}\right] \text{.} 
\label{approxFloatPhi}
\end{equation}
Thus if $T_{iMPSE}\approx T_{eMPSE}$ and $\gamma_{se}\ll1$,
\begin{equation}
\frac{e\left( \Phi_{W}^{f}-\Phi_{MPSE}^{f}\right) }{k_{B}T_{eMPSE}}\approx%
\frac{1}{2}\ln\left[ \frac{4\pi m_{e}}{m_{i}}\right] \text{.} 
\label{simpleFloatPhi}
\end{equation}

\bigskip

The appropriate choice of the free input parameters is $C_{hi2}=[2+\frac{%
e\left( \Phi_{MPSE}-\Phi_{W}\right) }{k_{B}T_{eMPSE}}]\left( \frac{\delta n}{%
\left\langle n\right\rangle }\right) _{MPSE}$, $C_{hi3}=\frac{7}{2}+\frac{1}{%
\sqrt{\pi}}+\frac{e\left( \Phi_{MPSE}-\Phi _{W}\right) }{k_{B}T_{eMPSE}}$
and all other $C_{hij}=0$ for the ion heat source and $C_{he3}=\sqrt{\frac{2%
}{\pi}}\exp\left[ \frac{e\left( \Phi _{W}-\Phi_{MPSE}\right) }{k_{B}T_{eMPSE}%
}\right] $, $C_{he4}=2\left( \frac{\delta n}{\left\langle n\right\rangle }%
\right) _{MPSE}$\ and all other $C_{hej}=0$ for the electron heat source.

\bigskip

\subsection{Standard Electric Charge Sources}

\bigskip

The electric charge source in each boundary cell is given by equation:
\begin{equation}
S_{ch}=-j_{d}S_{0}=J_{ex}-J_{ix}\text{,}   \label{chargeS}
\end{equation}
where $j_{d}$ is given by equation (\ref{jdTotal}) and $S_{0}$ is given by
equation (\ref{6}). In the general case, when $j_{d}\neq0$, the appropriate
choice of the free input parameters is $C_{ch6}=1$, $C_{ch7}=1-\gamma_{se}$
and all other $C_{chj}=0$. The non-zero numerical values of the electric
potential at the divertor plates, $\Phi_{d}^{\left( 0\right) }$, due to the
electric bias, can be specified in the input file \texttt{b2ai.dat}, which
describes the initial plasma state ($\Phi_{d}^{\left( 0\right) }=0$ is a
default value). In a case of electrically floating divertor plates, $j_{d}=0$%
, so that $S_{ch}=0$ in each boundary cell. It can be obtained by choosing
all $C_{chj}=0$.

\bigskip

\subsection{Conclusion}

\bigskip

It is obvious from this relatively simple example that only the expressions
for the particle sources, i.e., the numerical implementation of the boundary
conditions for the continuity equation, can adequately reproduce the
physical model of the magnetized boundary plasma, after some minor
corrections. It is necessary to include in the source code some more
complicated non-linear dependences of the boundary sources on the
macroscopic plasma parameters, so that the boundary conditions for the
higher moments' fluid equations can be correctly applied. The approximate
fluid model used in the upper example, which is based on the assumption that
the ion velocity distribution at the magnetic presheath edge is
drifting-Maxwellian distribution, can be a good starting point in the effort
to improve the standard expressions for the boundary sources.

\bigskip

\section{Boundary Sources Written by D. P. Coster\protect\bigskip}

\bigskip

\subsection{Introduction}

\bigskip

In order to improve the numerical implementation of the boundary conditions,
an alternative method for calculating the boundary sources has been
developed by D. P. Coster from IPP Garching. Their dependence on the plasma
parameters calculated with the bulk plasma equations is mostly the same as
for the standard boundary sources, but the way how the free input parameters
are specified and used in the calculations of the coefficients $S_{pj\alpha}$%
, $S_{mj\alpha}$, $S_{hij}$, $S_{hej}$ and $S_{ch}$ is different. The B2.5
user can choose between the standard (default) version and this version of
the boundary sources by specifying the numerical value of the input
parameter \texttt{boundary\_namelist} in the \texttt{b2mn.dat }input file
(default is zero). In order to use the non-standard version of the boundary
sources, a line of text: \texttt{%
%TCIMACRO{\UNICODE{0xb4}}%
%BeginExpansion
\'{}%
%EndExpansion
b2stbc\_boundary\_namelist%
%TCIMACRO{\UNICODE{0xb4}}%
%BeginExpansion
\'{}%
%EndExpansion
%TCIMACRO{\UNICODE{0xb4}}%
%BeginExpansion
\'{}%
%EndExpansion
1%
%TCIMACRO{\UNICODE{0xb4}}%
%BeginExpansion
\'{}%
%EndExpansion
}, should be added into the \texttt{b2mn.dat} file. The exact specifications
of the types of the boundary conditions and numerical values of the free
input parameters, which are used in the calculations of the non-standard
boundary sources, should be written in the \texttt{b2.boundary.parameters}
file. Following specifications of the boundary sources and corresponding
labels should be written in the \texttt{b2.boundary. parameters} file:

\begin{itemize}
\item  \texttt{\&BOUNDARY} (the first line or beginning of the file)

\item  \texttt{NBC=6,} (defines the number of the boundary conditions or the
number of the boundary regions)

\item  \texttt{BCCHAR= 
%TCIMACRO{\UNICODE{0xb4}}%
%BeginExpansion
\'{}%
%EndExpansion
S 
%TCIMACRO{\UNICODE{0xb4}}%
%BeginExpansion
\'{}%
%EndExpansion
, 
%TCIMACRO{\UNICODE{0xb4}}%
%BeginExpansion
\'{}%
%EndExpansion
S%
%TCIMACRO{\UNICODE{0xb4}}%
%BeginExpansion
\'{}%
%EndExpansion
, 
%TCIMACRO{\UNICODE{0xb4}}%
%BeginExpansion
\'{}%
%EndExpansion
S%
%TCIMACRO{\UNICODE{0xb4}}%
%BeginExpansion
\'{}%
%EndExpansion
, 
%TCIMACRO{\UNICODE{0xb4}}%
%BeginExpansion
\'{}%
%EndExpansion
N%
%TCIMACRO{\UNICODE{0xb4}}%
%BeginExpansion
\'{}%
%EndExpansion
, 
%TCIMACRO{\UNICODE{0xb4}}%
%BeginExpansion
\'{}%
%EndExpansion
W%
%TCIMACRO{\UNICODE{0xb4}}%
%BeginExpansion
\'{}%
%EndExpansion
, 
%TCIMACRO{\UNICODE{0xb4}}%
%BeginExpansion
\'{}%
%EndExpansion
E%
%TCIMACRO{\UNICODE{0xb4}}%
%BeginExpansion
\'{}%
%EndExpansion
,} (\texttt{BCCHAR(1...NBC)} defines the boundary as \texttt{%
%TCIMACRO{\UNICODE{0xb4}}%
%BeginExpansion
\'{}%
%EndExpansion
S%
%TCIMACRO{\UNICODE{0xb4}}%
%BeginExpansion
\'{}%
%EndExpansion
, 
%TCIMACRO{\UNICODE{0xb4}}%
%BeginExpansion
\'{}%
%EndExpansion
N%
%TCIMACRO{\UNICODE{0xb4}}%
%BeginExpansion
\'{}%
%EndExpansion
, 
%TCIMACRO{\UNICODE{0xb4}}%
%BeginExpansion
\'{}%
%EndExpansion
W%
%TCIMACRO{\UNICODE{0xb4}}%
%BeginExpansion
\'{}%
%EndExpansion
, 
%TCIMACRO{\UNICODE{0xb4}}%
%BeginExpansion
\'{}%
%EndExpansion
E%
%TCIMACRO{\UNICODE{0xb4}}%
%BeginExpansion
\'{}%
%EndExpansion
})

\item  \texttt{BCPOS=...,} (\texttt{BCPOS(1...NBC)} defines the position of
the boundary either in ix (\texttt{%
%TCIMACRO{\UNICODE{0xb4}}%
%BeginExpansion
\'{}%
%EndExpansion
W%
%TCIMACRO{\UNICODE{0xb4}}%
%BeginExpansion
\'{}%
%EndExpansion
} or \texttt{%
%TCIMACRO{\UNICODE{0xb4}}%
%BeginExpansion
\'{}%
%EndExpansion
E%
%TCIMACRO{\UNICODE{0xb4}}%
%BeginExpansion
\'{}%
%EndExpansion
}) or iy (\texttt{%
%TCIMACRO{\UNICODE{0xb4}}%
%BeginExpansion
\'{}%
%EndExpansion
S%
%TCIMACRO{\UNICODE{0xb4}}%
%BeginExpansion
\'{}%
%EndExpansion
}or \texttt{%
%TCIMACRO{\UNICODE{0xb4}}%
%BeginExpansion
\'{}%
%EndExpansion
N%
%TCIMACRO{\UNICODE{0xb4}}%
%BeginExpansion
\'{}%
%EndExpansion
}))

\item  \texttt{BCSTART=...,} (\texttt{BCSTART(1...NBC)} defines the start of
the boundary region either in ix (\texttt{%
%TCIMACRO{\UNICODE{0xb4}}%
%BeginExpansion
\'{}%
%EndExpansion
S%
%TCIMACRO{\UNICODE{0xb4}}%
%BeginExpansion
\'{}%
%EndExpansion
}or \texttt{%
%TCIMACRO{\UNICODE{0xb4}}%
%BeginExpansion
\'{}%
%EndExpansion
N%
%TCIMACRO{\UNICODE{0xb4}}%
%BeginExpansion
\'{}%
%EndExpansion
}) or iy (\texttt{%
%TCIMACRO{\UNICODE{0xb4}}%
%BeginExpansion
\'{}%
%EndExpansion
W%
%TCIMACRO{\UNICODE{0xb4}}%
%BeginExpansion
\'{}%
%EndExpansion
}or \texttt{%
%TCIMACRO{\UNICODE{0xb4}}%
%BeginExpansion
\'{}%
%EndExpansion
E%
%TCIMACRO{\UNICODE{0xb4}}%
%BeginExpansion
\'{}%
%EndExpansion
}))

\item  \texttt{BCEND=...,} (\texttt{BCEND(1...NBC)} defines the end of the
boundary region either in ix (\texttt{%
%TCIMACRO{\UNICODE{0xb4}}%
%BeginExpansion
\'{}%
%EndExpansion
S%
%TCIMACRO{\UNICODE{0xb4}}%
%BeginExpansion
\'{}%
%EndExpansion
}or \texttt{%
%TCIMACRO{\UNICODE{0xb4}}%
%BeginExpansion
\'{}%
%EndExpansion
N%
%TCIMACRO{\UNICODE{0xb4}}%
%BeginExpansion
\'{}%
%EndExpansion
}) or iy (\texttt{%
%TCIMACRO{\UNICODE{0xb4}}%
%BeginExpansion
\'{}%
%EndExpansion
W%
%TCIMACRO{\UNICODE{0xb4}}%
%BeginExpansion
\'{}%
%EndExpansion
}or \texttt{%
%TCIMACRO{\UNICODE{0xb4}}%
%BeginExpansion
\'{}%
%EndExpansion
E%
%TCIMACRO{\UNICODE{0xb4}}%
%BeginExpansion
\'{}%
%EndExpansion
}))

\item  \texttt{BCENE=...,} (\texttt{BCENE(1...NBC)=1-10} defines the type of
the boundary condition for the electron energy equation)

\item  \texttt{BCENI=...,} (\texttt{BCENI(1...NBC)=1-10} defines the type of
the boundary condition for the ion energy equation)

\item  \texttt{BCPOT=...,} (\texttt{BCPOT(1...NBC)=1-9} defines the type of
the boundary condition for the electric potential equation)

\item  \texttt{ENEPAR=...,} (\texttt{ENEPAR(1...NBC, 1)} and/or \texttt{%
ENEPAR(1...NBC, 2)} are free input parameters for the boundary condition for
the electron energy equation, chosen according to the corresponding \texttt{%
BCENE(1...NBC)})

\item  \texttt{ENIPAR=...,} (\texttt{ENIPAR(1...NBC, 1)} and/or \texttt{%
ENIPAR(1...NBC, 2)} are free input parameters for the boundary condition for
the ion energy equation, chosen according to the corresponding \texttt{%
BCENI(1...NBC)})

\item  \texttt{POTPAR=...,} (\texttt{POTPAR(1...NBC, 1)} and/or \texttt{%
POTPAR(1...NBC, 2)} are free input parameters for the boundary condition for
the electric potential equation, chosen according to the corresponding 
\texttt{BCPOT(1...NBC)})

\item  \texttt{BCCON=...,} (\texttt{BCCON(1...NS, 1...NBC)=1-10} defines the
type of the boundary condition for the continuity equation for species 
\texttt{IS=1...NS})

\item  \texttt{CONPAR=...,} (\texttt{CONPAR(1...NS, 1...NBC, 1)} and/or 
\texttt{CONPAR(1...NS, 1...NBC, 2)} are free input parameters for the
boundary condition for the continuity equation for species \texttt{IS=1...NS}%
, chosen according to the corresponding \texttt{BCCON(1...NS, 1...NBC)})

\item  \texttt{BCMOM=...,} (\texttt{BCMOM(1...NS, 1...NBC)=1-11} defines the
type of the boundary condition for the momentum equation for species \texttt{%
IS=1...NS})

\item  \texttt{MOMPAR=...,} (\texttt{MOMPAR(1...NS, 1...NBC, 1)} and/or 
\texttt{MOMPAR(1...NS, 1...NBC, 2)} are free input parameters for the
boundary condition for the momentum equation for species \texttt{IS=1...NS},
chosen according to the corresponding \texttt{BCMOM(1...NS, 1...NBC)})

\item  \texttt{GAMMAI=..., GAMMAE=...,} (ion adiabatic coefficient,
secondary electron emission coefficient)

\item  other labels for various input options, which are not implemented yet
(e.g., \texttt{LBNDUSR=t} (true) calls the user-routine \texttt{bndusr.f},
otherwise \texttt{LBNDUSR=f} (false)),

\item  \texttt{/} (the last line or end of the file).
\end{itemize}

\bigskip

If the value of \texttt{boundary\_namelist} is zero or the upper line of
text is not added, the standard boundary sources will be calculated.

\bigskip

This report describes the non-standard boundary sources included in the
version of \texttt{b2stbc.f} file from September 20, 1999 (see also [\ref
{mcoster}]). Some possible input parameter options for the non-standard
boundary sources have not been implemented yet. There are also several
certain or possible programming errors in the source code, in addition to
some theoretically questionable and/or insufficiently justified
calculations. These possible errors are also included in the analytical
expressions presented in this report. Since this part of the B2.5 code is
still not completely developed and debugged, these insufficiencies are not
analysed in detail in this report. However, it can be concluded that these
new boundary sources do not include significant qualitative improvements of
the standard boundary sources, or in some cases they can even less
accurately describe the boundary conditions compared with the standard
boundary sources.

\bigskip

\subsection{Analytical Expressions for the ''New'' Boundary Sources}

\bigskip

\subsubsection{Physical Parameters and Numerical Constants}

\bigskip

Several new numerical constants, which are related to the physical
velocities and particle density are defined. Their numerical values are the
following:

\begin{equation}
EPSLON=10^{-10}\equiv\mathcal{\epsilon}\text{,}   \label{c1}
\end{equation}

\begin{equation}
NE00=10^{20}\equiv\mathcal{N}_{e}\text{,}   \label{c2}
\end{equation}

\begin{equation}
UP00=\sqrt{\frac{100eV}{m_{p}}}\circeq97873\frac{m}{s}\equiv\mathcal{U}_{p}%
\text{,}   \label{c3}
\end{equation}

\begin{equation}
UBIG=\frac{UP00}{EPSLON}\equiv\mathcal{U}   \label{c4}
\end{equation}
and

\begin{equation}
VBIG=10^{-2}UBIG\equiv\mathcal{V}\text{.}   \label{c5}
\end{equation}
Their physical background is not explained in the comments. Anomalous
diffusion coefficients $D_{an}^{\alpha}$ and $D_{ap}^{\alpha}$, viscosities $%
\eta_{a}^{\alpha}$, heat conductivities $\kappa_{a}^{e}$ and $\kappa_{a}^{i}$
and electrical conductivity $\sigma_{a}$ are also used in the calculations
of the boundary sources. $s_{x}$ is the area of the face between the cells
(ix, iy) and (ix+1,iy) and $s_{y}$ between the cells (ix,iy) and (ix,iy+1).
In the expressions for the boundary sources they are calculated at the
boundary cells. $u_{\alpha}$ is the poloidal component of the drift velocity
of the atomic species $\alpha$, $u_{D\alpha}$ is the poloidal component of
the $\mathbf{E\times B}$ drift velocity and $b_{x}=\frac{B_{x}}{B}$ is the
magnetic field pitch. The ion sound velocity for the multispecies plasma is
defined as:

\begin{equation}
c_{s}=\sqrt{\frac{p}{\rho}}\text{,}   \label{c6}
\end{equation}
where $p$ is the total pressure and $\rho=m_{p}\sum_{\alpha}A_{\alpha
}n_{\alpha}$ is the total mass density. The ion sound velocity of the
species $\alpha$ is defined as:

\begin{equation}
c_{s\alpha}=\sqrt{\frac{\gamma_{i}\widehat{T}_{i}+Z_{\alpha}^{2}r_{\alpha }%
\widehat{T}_{e}}{m_{p}A_{\alpha}}}\text{,}   \label{c7}
\end{equation}
where $\gamma_{i}=\frac{c_{pi}}{c_{vi}}$ is the ratio of the ion specific
heats or the ion adiabatic coefficient ($\gamma_{i}=\frac{5}{3}$ for the
adiabatic fluid model, $\gamma_{i}=1$ for the isothermal fluid model, $%
\gamma_{i}=3$ for the one-dimensional kinetic model [\ref{mriemann91}]) and $%
r_{\alpha}=\frac{n_{\alpha}}{n_{e}}$. $C_{leak}$ is a dimensionless leakage
coefficient. $\gamma_{se}$ is the secondary electron emission coefficient.
The other physical parameters are described in the following text, when they
are used in the expressions for the boundary sources.

\bigskip

\subsubsection{Boundary Conditions for Continuity Equations}

\bigskip

\paragraph{\texttt{BCCON=1} ... Prescribe the Value of the Density}

\bigskip

\begin{equation*}
\mathtt{CONPAR(IS,IB,1)}\equiv N_{\alpha}
\end{equation*}

\bigskip

\begin{itemize}
\item  SOUTH and NORTH boundary:
\end{itemize}

\begin{equation}
S_{p0\alpha}=\mathcal{V}N_{\alpha}s_{y}\text{,}   \label{c8}
\end{equation}

\begin{equation}
S_{p1\alpha}=-\mathcal{V}s_{y}\text{,}   \label{c9}
\end{equation}

\begin{equation}
S_{p\alpha}=S_{p0\alpha}+S_{p1\alpha}n_{\alpha}=\mathcal{V}s_{y}\left(
N_{\alpha}-n_{\alpha}\right) \text{.}   \label{c10}
\end{equation}

\begin{itemize}
\item  WEST and EAST boundary:
\end{itemize}

\begin{equation}
S_{p0\alpha}=\mathcal{U}N_{\alpha}s_{x}\text{,}   \label{c11}
\end{equation}

\begin{equation}
S_{p1\alpha}=-\mathcal{U}s_{x}\text{,}   \label{c12}
\end{equation}

\begin{equation}
S_{p\alpha}=S_{p0\alpha}+S_{p1\alpha}n_{\alpha}=\mathcal{U}s_{x}\left(
N_{\alpha}-n_{\alpha}\right) \text{.}   \label{c13}
\end{equation}

\bigskip

\paragraph{\texttt{BCCON=2} ... Prescribe the Value of the Gradient of the
Density}

\bigskip

\begin{equation*}
\mathtt{CONPAR(IS,IB,1)}\equiv\partial_{x}N_{\alpha},\text{ }\partial
_{y}N_{\alpha}
\end{equation*}

\bigskip

\begin{itemize}
\item  SOUTH boundary:
\end{itemize}

\begin{equation}
l_{y}=\frac{1}{2}\left( h_{y1(S)}+h_{y1(S+1)}\right) \text{,}   \label{c14}
\end{equation}

\begin{equation}
S_{p0\alpha}=\mathcal{V}s_{y}\left[ \left( -\partial_{y}N_{\alpha}\right)
l_{y}+n_{\alpha}\right] \text{,}   \label{c15}
\end{equation}

\begin{equation}
S_{p1\alpha}=-\mathcal{V}s_{y}\text{,}   \label{c16}
\end{equation}

\begin{equation}
S_{p\alpha}=S_{p0\alpha}+S_{p1\alpha}n_{\alpha}=\mathcal{V}s_{y}\left(
-\partial_{y}N_{\alpha}\right) l_{y}\text{.}   \label{c17}
\end{equation}

\begin{itemize}
\item  NORTH boundary:
\end{itemize}

\begin{equation}
l_{y}=\frac{1}{2}\left( h_{y1(N-1)}+h_{y1(N)}\right) \text{,}   \label{c18}
\end{equation}

\begin{equation}
S_{p0\alpha}=\mathcal{V}s_{y}\left[ \left( \partial_{y}N_{\alpha}\right)
l_{y}+n_{\alpha}\right] \text{,}   \label{c19}
\end{equation}

\begin{equation}
S_{p1\alpha}=-\mathcal{V}s_{y}\text{,}   \label{c20}
\end{equation}

\begin{equation}
S_{p\alpha}=S_{p0\alpha}+S_{p1\alpha}n_{\alpha}=\mathcal{V}s_{y}\left(
\partial_{y}N_{\alpha}\right) l_{y}\text{.}   \label{c21}
\end{equation}

\begin{itemize}
\item  WEST boundary:
\end{itemize}

\begin{equation}
l_{x}=\frac{1}{2}\left( h_{x(W)}+h_{x(W+1)}\right) \text{,}   \label{c22}
\end{equation}

\begin{equation}
S_{p0\alpha}=\mathcal{U}s_{x}\left[ \left( -\partial_{x}N_{\alpha}\right)
l_{x}+n_{\alpha}\right] \text{,}   \label{c23}
\end{equation}

\begin{equation}
S_{p1\alpha}=-\mathcal{U}s_{x}\text{,}   \label{c24}
\end{equation}

\begin{equation}
S_{p\alpha}=S_{p0\alpha}+S_{p1\alpha}n_{\alpha}=\mathcal{U}s_{x}\left(
-\partial_{x}N_{\alpha}\right) l_{x}\text{.}   \label{c25}
\end{equation}

\begin{itemize}
\item  EAST boundary:
\end{itemize}

\begin{equation}
l_{x}=\frac{1}{2}\left( h_{x(E-1)}+h_{x(E)}\right) \text{,}   \label{c26}
\end{equation}

\begin{equation}
S_{p0\alpha}=\mathcal{U}s_{x}\left[ \left( \partial_{x}N_{\alpha}\right)
l_{x}+n_{\alpha}\right] \text{,}   \label{c27}
\end{equation}

\begin{equation}
S_{p1\alpha}=-\mathcal{U}s_{x}\text{,}   \label{c28}
\end{equation}

\begin{equation}
S_{p\alpha}=S_{p0\alpha}+S_{p1\alpha}n_{\alpha}=\mathcal{U}s_{x}\left(
\partial_{x}N_{\alpha}\right) l_{x}\text{.}   \label{c29}
\end{equation}

\bigskip

\paragraph{\texttt{BCCON=3} ... Sheath Conditions (Only for the WEST and
EAST Boundary)}

\bigskip

\begin{equation}
S_{p0\alpha}=\mathcal{U}n_{\alpha}s_{x},   \label{ccc30}
\end{equation}

\begin{equation}
S_{p1\alpha}=-\mathcal{U}s_{x}\text{,}   \label{c31}
\end{equation}

\begin{equation}
S_{p\alpha}=S_{p0\alpha}+S_{p1\alpha}n_{\alpha}=0\text{.}^{\left(
\ast\right) }   \label{c32}
\end{equation}

$\bigskip$

$^{\left( \ast\right) }$WARNING: PROBABLY ERROR - $S_{p1\alpha}$ should have
some other value, so that $S_{p\alpha}\neq0$!

\bigskip

\paragraph{\texttt{BCCON=4} ... Prescribe the Value of the Density
''Weakly'', Using as the Control Parameter Some Prescribed Velocity}

\bigskip

\begin{equation*}
\mathtt{CONPAR(IS,IB,1)}\equiv N_{\alpha}\text{,}
\end{equation*}

\begin{equation*}
\mathtt{CONPAR(IS,IB,2)}\equiv U_{\alpha}
\end{equation*}

\bigskip

\begin{itemize}
\item  for all boundaries ($N_{\alpha}$ and $U_{\alpha}$ are different at
each boundary):
\end{itemize}

\begin{equation}
S_{p0\alpha}=N_{\alpha}U_{\alpha}s\text{,}   \label{c33}
\end{equation}

\begin{equation}
S_{p1\alpha}=-U_{\alpha}s\text{,}   \label{c34}
\end{equation}

\begin{equation}
S_{p\alpha}=S_{p0\alpha}+S_{p1\alpha}n_{\alpha}=U_{\alpha}s\left( N_{\alpha
}-n_{\alpha}\right) \text{.}   \label{c35}
\end{equation}

\bigskip

\paragraph{\texttt{BCCON=5} ... Prescribe the Value of the Particle Flux per
Unit Area}

\bigskip

\begin{equation*}
\mathtt{CONPAR(IS,IB,1)}\equiv\Gamma_{\alpha x},\text{ }\Gamma_{\alpha y}
\end{equation*}

\bigskip

\begin{itemize}
\item  SOUTH and NORTH boundary:
\end{itemize}

\begin{equation}
S_{p0\alpha}=\Gamma_{\alpha y}s_{y}\text{,}   \label{c36}
\end{equation}

\begin{equation}
S_{p1\alpha}=0\text{,}   \label{c37}
\end{equation}

\begin{equation}
S_{p\alpha}=S_{p0\alpha}+S_{p1\alpha}n_{\alpha}=\Gamma_{\alpha y}s_{y}\text{.%
}   \label{c38}
\end{equation}

\begin{itemize}
\item  WEST and EAST boundary:
\end{itemize}

\begin{equation}
S_{p0\alpha}=\Gamma_{\alpha x}s_{x}\text{,}   \label{c39}
\end{equation}

\begin{equation}
S_{p1\alpha}=0\text{,}   \label{c40}
\end{equation}

\begin{equation}
S_{p\alpha}=S_{p0\alpha}+S_{p1\alpha}n_{\alpha}=\Gamma_{\alpha x}s_{x}\text{.%
}   \label{c41}
\end{equation}

\bigskip

\paragraph{\texttt{BCCON=6} ... Prescribe the Value of the Total Particle
Flux for a Constant Density (Only for the NORTH and SOUTH Boundary)}

\bigskip

\begin{equation*}
\mathtt{CONPAR(IS,IB,1)}\equiv\mathcal{S}_{n}
\end{equation*}

\bigskip

\begin{itemize}
\item  NORTH boundary:
\end{itemize}

\begin{equation}
s_{1}=\sum_{j}^{N}\frac{s_{yj}}{\frac{1}{2}\left( h_{y1\left( j,N-1\right)
}+h_{y1\left( j,N\right) }\right) }\text{,}   \label{c42}
\end{equation}

\begin{equation}
s_{2}=\sum_{j}^{N}\frac{n_{\alpha j}s_{yj}}{\frac{1}{2}\left( h_{y1\left(
j,N-1\right) }+h_{y1\left( j,N\right) }\right) }\text{,}   \label{c43}
\end{equation}

\begin{itemize}
\item  SOUTH boundary:
\end{itemize}

\begin{equation}
s_{1}=\sum_{j}^{S}\frac{s_{yj}}{\frac{1}{2}\left( h_{y1\left( j,S\right)
}+h_{y1\left( j,S+1\right) }\right) }\text{,}   \label{c44}
\end{equation}

\begin{equation}
s_{2}=\sum_{j}^{S}\frac{n_{\alpha j}s_{yj}}{\frac{1}{2}\left( h_{y1\left(
j,S\right) }+h_{y1\left( j,S+1\right) }\right) },   \label{cs45}
\end{equation}
where $\sum_{j}^{N}$ and $\sum_{j}^{S}$ are summations over the boundary
cells at the NORTH and SOUTH boundary respectively. The particle sources are
defined as:

\begin{equation}
S_{p0\alpha}=\mathcal{V}s_{y}\left( \mathcal{S}_{n}/D_{an}+s_{2}\right)
/s_{1}\text{,}   \label{c46}
\end{equation}

\begin{equation}
S_{p1\alpha}=-\mathcal{V}s_{y}\text{,}   \label{c47}
\end{equation}

\begin{equation}
S_{p\alpha}=S_{p0\alpha}+S_{p1\alpha}n_{\alpha}=\mathcal{V}s_{y}\left[ 
\mathcal{S}_{n}/\left( D_{an}s_{1}\right) +s_{2}/s_{1}-n_{\alpha}\right] 
\text{.}   \label{c48}
\end{equation}

\bigskip

\paragraph{\texttt{BCCON=7} ... Prescribe the Value of the Density as a
Function of Other Plasma Parameters (Not Yet Available)}

\bigskip

\bigskip

\paragraph{\texttt{BCCON=8} ... Prescribe the Value of the Total Particle
Flux with a Constant Flux Density}

\bigskip

\begin{equation*}
\mathtt{CONPAR(IS,IB,1)}\equiv\mathcal{S}_{\Gamma_{x}},\text{ }\mathcal{S}%
_{\Gamma_{y}}
\end{equation*}

\bigskip

\begin{itemize}
\item  SOUTH and NORTH boundary:
\end{itemize}

\begin{equation}
s_{tot}=\sum_{j}^{S\text{ or }N}s_{yj}\text{,}   \label{c49}
\end{equation}

\begin{equation}
S_{p0\alpha}=\left( \mathcal{S}_{\Gamma_{y}}/s_{tot}\right) s_{y}\text{,} 
\label{c50}
\end{equation}

\begin{equation}
S_{p1\alpha}=0\text{,}   \label{c51}
\end{equation}

\begin{equation}
S_{p\alpha}=S_{p0\alpha}+S_{p1\alpha}n_{\alpha}=\left( \mathcal{S}%
_{\Gamma_{y}}/s_{tot}\right) s_{y}\text{.}   \label{c52}
\end{equation}

\begin{itemize}
\item  WEST and EAST boundary:
\end{itemize}

\begin{equation}
s_{tot}=\sum_{j}^{W\text{ or }E}s_{xj}\text{,}   \label{c53}
\end{equation}

\begin{equation}
S_{p0\alpha}=\left( \mathcal{S}_{\Gamma_{x}}/s_{tot}\right) s_{x}\text{,} 
\label{c54}
\end{equation}

\begin{equation}
S_{p1\alpha}=0\text{,}   \label{c55}
\end{equation}

\begin{equation}
S_{p\alpha}=S_{p0\alpha}+S_{p1\alpha}n_{\alpha}=\left( \mathcal{S}%
_{\Gamma_{x}}/s_{tot}\right) s_{x}\text{.}   \label{c56}
\end{equation}

\bigskip

\paragraph{\texttt{BCCON=9} ... Prescribe the Value of the Decay Length for
the Density (Only for the SOUTH and NORTH Boundary)}

\bigskip

\begin{equation*}
\mathtt{CONPAR(IS,IB,1)}\equiv
l_{n_{\alpha}}=n_{\alpha}/\partial_{y}n_{\alpha}
\end{equation*}

\bigskip

\begin{equation}
S_{p0\alpha}=0\text{,}   \label{c57}
\end{equation}

\begin{equation}
S_{p1\alpha}=-\left[ D_{an_{\alpha}}+D_{ap_{\alpha}}\left( Z_{\alpha }%
\widehat{T}_{e}+\widehat{T}_{i}\right) \right] s_{y}/l_{n_{\alpha}}\text{,} 
\label{c58}
\end{equation}

\begin{equation}
S_{p\alpha}=S_{p0\alpha}+S_{p1\alpha}n_{\alpha}=-\left[ D_{an_{%
\alpha}}+D_{ap_{\alpha}}\left( Z_{\alpha}\widehat{T}_{e}+\widehat{T}%
_{i}\right) \right] s_{y}n_{\alpha}/l_{n_{\alpha}}\text{.}   \label{c59}
\end{equation}

\bigskip

\paragraph{\texttt{BCCON=10} ... Prescribe the Leakage Coefficient (Leakage
Option for the Density, Only for the SOUTH and NORTH Boundary)}

\bigskip

\begin{equation*}
\mathtt{CONPAR(IS,IB,1)}\equiv C_{leak}
\end{equation*}

\bigskip

\begin{equation}
S_{p0\alpha}=0\text{,}   \label{c60}
\end{equation}

\begin{equation}
S_{p1\alpha}=\frac{V}{h_{y1}}\sqrt{\frac{Z_{\alpha}\widehat{T}_{e}+\widehat
{T}_{i}}{m_{p}A_{\alpha}}}C_{leak}\text{,}   \label{c70}
\end{equation}

\begin{equation}
S_{p\alpha}=S_{p0\alpha}+S_{p1\alpha}n_{\alpha}=\frac{V}{h_{y1}}\sqrt {\frac{%
Z_{\alpha}\widehat{T}_{e}+\widehat{T}_{i}}{m_{p}A_{\alpha}}%
}C_{leak}n_{\alpha}\text{.}   \label{c71}
\end{equation}

\begin{itemize}
\item  total pumping rate of the species $\alpha$ at the SOUTH and NORTH
boundary (summation over all boundary cells, added by H. B\"{u}rbaumer):
\end{itemize}

\begin{equation}
R_{p\alpha}=\sum_{j}^{S\text{ or }N}S_{p\alpha j}=\sum_{j}^{S\text{ or }%
N}S_{p1\alpha j}n_{\alpha j}\text{.}   \label{c72}
\end{equation}

\bigskip

\subsubsection{Boundary Conditions for Momentum Equations}

\bigskip

\paragraph{\texttt{BCMOM=1} ... Prescribe the Value of the Parallel Velocity}

\bigskip

\begin{equation*}
\mathtt{MOMPAR(IS,IB,1)}\equiv U_{\alpha}
\end{equation*}

\bigskip

\begin{itemize}
\item  SOUTH and NORTH boundary:
\end{itemize}

\begin{equation}
S_{m0\alpha}=\mathcal{VN}_{e}A_{\alpha}U_{\alpha}s_{y}\text{,}   \label{c73}
\end{equation}

\begin{equation}
S_{m1\alpha}=-\mathcal{VN}_{e}A_{\alpha}s_{y}\text{,}   \label{c74}
\end{equation}

\begin{equation}
S_{m\alpha}=S_{m0\alpha}+S_{m1\alpha}u_{\alpha}=\mathcal{VN}_{e}A_{\alpha
}s_{y}\left( U_{\alpha}-u_{\alpha}\right) \text{.}   \label{c75}
\end{equation}

\begin{itemize}
\item  WEST and EAST boundary:
\end{itemize}

\begin{equation}
S_{m0\alpha}=\mathcal{UN}_{e}A_{\alpha}U_{\alpha}s_{x}\text{,}   \label{c76}
\end{equation}

\begin{equation}
S_{m1\alpha}=-\mathcal{UN}_{e}A_{\alpha}s_{x}\text{,}   \label{c77}
\end{equation}

\begin{equation}
S_{m\alpha}=S_{m0\alpha}+S_{m1\alpha}u_{\alpha}=\mathcal{UN}_{e}A_{\alpha
}s_{x}\left( U_{\alpha}-u_{\alpha}\right) \text{.}   \label{c78}
\end{equation}

\bigskip

\paragraph{\texttt{BCMOM=2} ... Prescribe the Gradient of the Parallel
Velocity}

\bigskip

\begin{equation*}
\mathtt{MOMPAR(IS,IB,1)}\equiv\partial_{x}U_{\alpha},\text{ }\partial
_{y}U_{\alpha}
\end{equation*}

\bigskip

\begin{itemize}
\item  NORTH boundary:
\end{itemize}

\begin{equation}
\overline{n}_{\alpha}=\frac{1}{2}\left( n_{\alpha\left( ix,N\right)
}+n_{\alpha\left( ix+1,N\right) }\right) \text{,}   \label{c79}
\end{equation}

\begin{equation}
l_{y}=\frac{1}{4}\left( h_{y1\left( ix,N-1\right) }+h_{y1\left(
ix+1,N-1\right) }\right) \text{,}   \label{c80}
\end{equation}

\begin{equation}
\widehat{s}=\frac{1}{2}\left( b_{x\left( ix,N-1\right) }+b_{x\left(
ix+1,N-1\right) }\right) \text{,}   \label{c81}
\end{equation}

\begin{equation}
\widehat{c}=\frac{1}{2}\left( \sqrt{1-b_{x\left( ix,N-1\right) }^{2}}+\sqrt{%
1-b_{x\left( ix+1,N-1\right) }^{2}}\right) \text{,}   \label{c82}
\end{equation}

\begin{equation}
S_{m0\alpha}=\mathcal{V}s_{y}A_{\alpha}\overline{n}_{\alpha}\left[ \left(
\partial_{y}U_{\alpha}\right) l_{y}+\left( u_{\alpha}-\widehat{c}u_{D\alpha
}\right) /\widehat{s}\right] \text{,}   \label{c83}
\end{equation}

\begin{equation}
S_{m1\alpha}=-\mathcal{V}s_{y}A_{\alpha}\overline{n}_{\alpha}\text{,} 
\label{c84}
\end{equation}

\begin{equation}
S_{m\alpha}=S_{m0\alpha}+S_{m1\alpha}u_{\alpha}=\mathcal{V}s_{y}A_{\alpha }%
\overline{n}_{\alpha}\left[ \left( \partial_{y}U_{\alpha}\right)
l_{y}+\left( u_{\alpha}-\widehat{c}u_{D\alpha}\right) /\widehat{s}-u_{\alpha}%
\right] \text{.}^{\left( \ast\right) }   \label{cc85}
\end{equation}

\bigskip

$^{\left( \ast\right) }$WARNING: PROBABLY\ ERROR - THE\ AUTHOR\ PROBABLY\
WANTED\ TO\ GET\ THE\ RESULT:

\begin{equation*}
S_{m\alpha}=\mathcal{V}s_{y}A_{\alpha}\overline{n}_{\alpha}\left(
\partial_{y}U_{\alpha}\right) l_{y}\text{.}
\end{equation*}

\bigskip

\begin{itemize}
\item  SOUTH boundary:
\end{itemize}

\begin{equation}
\overline{n}_{\alpha}=\frac{1}{2}\left( n_{\alpha\left( ix,S\right)
}+n_{\alpha\left( ix+1,S\right) }\right) \text{,}   \label{c86}
\end{equation}

\begin{equation}
l_{y}=\frac{1}{4}\left( h_{y1\left( ix,S+1\right) }+h_{y1\left(
ix+1,S+1\right) }\right) \text{,}   \label{c87}
\end{equation}

\begin{equation}
\widehat{s}=\frac{1}{2}\left( b_{x\left( ix,S+1\right) }+b_{x\left(
ix+1,S+1\right) }\right) \text{,}   \label{c88}
\end{equation}

\begin{equation}
\widehat{c}=\frac{1}{2}\left( \sqrt{1-b_{x\left( ix,S+1\right) }^{2}}+\sqrt{%
1-b_{x\left( ix+1,S+1\right) }^{2}}\right) \text{,}   \label{c89}
\end{equation}

\begin{equation}
S_{m0\alpha}=\mathcal{V}s_{y}A_{\alpha}\overline{n}_{\alpha}\left[ -\left(
\partial_{y}U_{\alpha}\right) l_{y}+\left( u_{\alpha}-\widehat{c}u_{D\alpha
}\right) /\widehat{s}\right] \text{,}   \label{c90}
\end{equation}

\begin{equation}
S_{m1\alpha}=-\mathcal{V}s_{y}A_{\alpha}\overline{n}_{\alpha}\text{,} 
\label{c91}
\end{equation}

\begin{equation}
S_{m\alpha}=S_{m0\alpha}+S_{m1\alpha}u_{\alpha}=\mathcal{V}s_{y}A_{\alpha }%
\overline{n}_{\alpha}\left[ -\left( \partial_{y}U_{\alpha}\right)
l_{y}+\left( u_{\alpha}-\widehat{c}u_{D\alpha}\right) /\widehat{s}-u_{\alpha}%
\right] \text{.}^{\left( \ast\right) }   \label{c92}
\end{equation}

\bigskip

$^{\left( \ast\right) }$WARNING: PROBABLY\ ERROR - THE\ AUTHOR\ PROBABLY\
WANTED\ TO\ GET\ THE\ RESULT:

\begin{equation*}
S_{m\alpha}=-\mathcal{V}s_{y}A_{\alpha}\overline{n}_{\alpha}\left(
\partial_{y}U_{\alpha}\right) l_{y}\text{.}
\end{equation*}

\bigskip

\begin{itemize}
\item  WEST boundary:
\end{itemize}

\begin{equation}
\overline{n}_{\alpha}=\frac{1}{2}\left( n_{\alpha\left( W,iy\right)
}+n_{\alpha\left( W+1,iy\right) }\right) \text{,}   \label{c93}
\end{equation}

\begin{equation}
l_{x}=\frac{1}{4}\left( h_{x\left( W,iy\right) }+h_{x\left( W+1,iy\right)
}\right) \text{,}   \label{c94}
\end{equation}

\begin{equation}
\widehat{s}=\frac{1}{2}\left( b_{x\left( W,iy\right) }+b_{x\left(
W+1,iy\right) }\right) \text{,}   \label{c95}
\end{equation}

\begin{equation}
\widehat{c}=\frac{1}{2}\left( \sqrt{1-b_{x\left( W,iy\right) }^{2}}+\sqrt{%
1-b_{x\left( W+1,iy\right) }^{2}}\right) \text{,}   \label{c96}
\end{equation}

\begin{equation}
S_{m0\alpha}=\mathcal{V}s_{x}A_{\alpha}\overline{n}_{\alpha}\left[ -\left(
\partial_{x}U_{\alpha}\right) l_{x}+\left( u_{\alpha}-\widehat{c}u_{D\alpha
}\right) /\widehat{s}\right] \text{,}   \label{c97}
\end{equation}

\begin{equation}
S_{m1\alpha}=-\mathcal{V}s_{x}A_{\alpha}\overline{n}_{\alpha}\text{,} 
\label{c98}
\end{equation}

\begin{equation}
S_{m\alpha}=S_{m0\alpha}+S_{m1\alpha}u_{\alpha}=\mathcal{V}s_{x}A_{\alpha }%
\overline{n}_{\alpha}\left[ -\left( \partial_{x}U_{\alpha}\right)
l_{x}+\left( u_{\alpha}-\widehat{c}u_{D\alpha}\right) /\widehat{s}-u_{\alpha}%
\right] \text{.}^{\left( \ast\right) }   \label{c99}
\end{equation}

\bigskip

$^{\left( \ast\right) }$WARNING: PROBABLY\ ERROR - THE\ AUTHOR\ PROBABLY\
WANTED\ TO\ GET\ THE\ RESULT:

\begin{equation*}
S_{m\alpha}=-\mathcal{V}s_{x}A_{\alpha}\overline{n}_{\alpha}\left(
\partial_{x}U_{\alpha}\right) l_{x}\text{.}
\end{equation*}

\bigskip

\begin{itemize}
\item  EAST boundary:
\end{itemize}

\begin{equation}
\overline{n}_{\alpha}=\frac{1}{2}\left( n_{\alpha\left( E-1,iy\right)
}+n_{\alpha\left( E,iy\right) }\right) \text{,}   \label{c100}
\end{equation}

\begin{equation}
l_{x}=\frac{1}{4}\left( h_{x\left( E-1,iy\right) }+h_{x\left( E,iy\right)
}\right) \text{,}   \label{c101}
\end{equation}

\begin{equation}
\widehat{s}=\frac{1}{2}\left( b_{x\left( E-1,iy\right) }+b_{x\left(
E,iy\right) }\right) \text{,}   \label{c102}
\end{equation}

\begin{equation}
\widehat{c}=\frac{1}{2}\left( \sqrt{1-b_{x\left( E-1,iy\right) }^{2}}+\sqrt{%
1-b_{x\left( E,iy\right) }^{2}}\right) \text{,}   \label{c103}
\end{equation}

\begin{equation}
S_{m0\alpha}=\mathcal{V}s_{x}A_{\alpha}\overline{n}_{\alpha}\left[ \left(
\partial_{x}U_{\alpha}\right) l_{x}+\left( u_{\alpha}-\widehat{c}u_{D\alpha
}\right) /\widehat{s}\right] \text{,}   \label{c104}
\end{equation}

\begin{equation}
S_{m1\alpha}=-\mathcal{V}s_{x}A_{\alpha}\overline{n}_{\alpha}\text{,} 
\label{c105}
\end{equation}

\begin{equation}
S_{m\alpha}=S_{m0\alpha}+S_{m1\alpha}u_{\alpha}=\mathcal{V}s_{x}A_{\alpha }%
\overline{n}_{\alpha}\left[ \left( \partial_{x}U_{\alpha}\right)
l_{x}+\left( u_{\alpha}-\widehat{c}u_{D\alpha}\right) /\widehat{s}-u_{\alpha}%
\right] \text{.}^{\left( \ast\right) }   \label{c106}
\end{equation}

\bigskip

$^{\left( \ast\right) }$WARNING: PROBABLY\ ERROR - THE\ AUTHOR\ PROBABLY\
WANTED\ TO\ GET\ THE\ RESULT:

\begin{equation*}
S_{m\alpha}=\mathcal{V}s_{x}A_{\alpha}\overline{n}_{\alpha}\left(
\partial_{x}U_{\alpha}\right) l_{x}\text{.}
\end{equation*}

\bigskip

\paragraph{\texttt{BCMOM=3} ... Sheath Conditions (Only for the WEST and
EAST Boundary), Mach Number as Input}

\bigskip

\begin{equation*}
\mathtt{MOMPAR(IS,IB,1)}\equiv\mathcal{M}_{\alpha1}, 
\end{equation*}

\begin{equation*}
\mathtt{MOMPAR(IS,IB,2)}\equiv\mathcal{M}_{\alpha2}
\end{equation*}

\bigskip

\begin{itemize}
\item  WEST boundary:
\end{itemize}

\bigskip

IF $\left( \mathcal{M}_{\alpha2}<0.5\right) $ THEN:

\begin{equation}
u_{out}=b_{x}\mathcal{M}_{11}c_{s}\text{,}   \label{c107}
\end{equation}

\begin{equation}
u_{p\alpha}=\frac{u_{out}}{b_{x}}+u_{D1}\text{,}   \label{c108}
\end{equation}

\bigskip

ELSE $\left( \mathcal{M}_{\alpha2}\geq0.5\right) $:

\begin{equation}
u_{test\alpha}=u_{\alpha\left( W+2,iy\right) }\left( 1+\frac{h_{x\left(
W+1,iy\right) }}{h_{x\left( W+2,iy\right) }}\right) -u_{\alpha\left(
W+3,iy\right) }\frac{h_{x\left( W+1,iy\right) }}{h_{x\left( W+2,iy\right) }}%
\text{,}   \label{c109}
\end{equation}

\begin{equation}
u_{out\alpha}=\max\left( \mathcal{M}_{\alpha1}c_{s\alpha}\text{, }%
-u_{test\alpha}\right) \text{,}   \label{c110}
\end{equation}

\begin{equation}
u_{p\alpha}=u_{out\alpha}+u_{D1}\text{,}   \label{c111}
\end{equation}
where $u_{D1}$ is the poloidal component of the $\mathbf{E\times B}$ drift
velocity.

\bigskip

\begin{equation}
S_{m0\alpha}=\mathcal{U}A_{\alpha}m_{p}n_{\alpha}s_{x}\left( -u_{p\alpha
}\right) \text{,}   \label{c112}
\end{equation}

\begin{equation}
S_{m1\alpha}=-\mathcal{U}A_{\alpha}m_{p}n_{\alpha}s_{x}\text{,} 
\label{c113}
\end{equation}

\begin{equation}
S_{m\alpha}=S_{m0\alpha}+S_{m1\alpha}u_{\alpha}=\mathcal{U}%
A_{\alpha}m_{p}n_{\alpha}s_{x}\left( -u_{p\alpha}-u_{\alpha}\right) \text{.} 
\label{c114}
\end{equation}

\bigskip

\begin{itemize}
\item  EAST boundary:
\end{itemize}

\bigskip

IF $\left( \mathcal{M}_{\alpha2}<0.5\right) $ THEN:

\begin{equation}
u_{out}=b_{x}\mathcal{M}_{11}c_{s}\text{,}   \label{c115}
\end{equation}

\begin{equation}
u_{p\alpha}=\frac{u_{out}}{b_{x}}-u_{D1}\text{,}   \label{c116}
\end{equation}

\bigskip

ELSE $\left( \mathcal{M}_{\alpha2}\geq0.5\right) $:

\begin{equation}
u_{test\alpha}=u_{\alpha\left( E-1,iy\right) }\left( 1+\frac{h_{x\left(
E-1,iy\right) }}{h_{x\left( E-2,iy\right) }}\right) -u_{\alpha\left(
E-2,iy\right) }\frac{h_{x\left( E-1,iy\right) }}{h_{x\left( E-2,iy\right) }}%
\text{,}   \label{c117}
\end{equation}

\begin{equation}
u_{out\alpha}=\max\left( \mathcal{M}_{\alpha1}c_{s\alpha}\text{, }%
u_{test\alpha}\right) \text{,}   \label{c118}
\end{equation}

\begin{equation}
u_{p\alpha}=u_{out\alpha}-u_{D1}\text{,}   \label{c119}
\end{equation}
where $u_{D1}$ is the poloidal component of the $\mathbf{E\times B}$ drift
velocity.

\bigskip

\begin{equation}
S_{m0\alpha}=\mathcal{U}A_{\alpha}m_{p}n_{\alpha}s_{x}u_{p\alpha}\text{,} 
\label{c120}
\end{equation}

\begin{equation}
S_{m1\alpha}=-\mathcal{U}A_{\alpha}m_{p}n_{\alpha}s_{x}\text{,} 
\label{c121}
\end{equation}

\begin{equation}
S_{m\alpha}=S_{m0\alpha}+S_{m1\alpha}u_{\alpha}=\mathcal{U}%
A_{\alpha}m_{p}n_{\alpha}s_{x}\left( u_{p\alpha}-u_{\alpha}\right) \text{.} 
\label{c122}
\end{equation}

\bigskip

\paragraph{\texttt{BCMOM=4} ... Prescribe the Value of the Velocity
''Weakly'', Using as the Control Parameter Some Prescribed Mass Flux}

\bigskip

\begin{equation*}
\mathtt{MOMPAR(IS,IB,1)}\equiv U_{\alpha}, 
\end{equation*}

\begin{equation*}
\mathtt{MOMPAR(IS,IB,2)}\equiv G_{\alpha x},\text{ }G_{\alpha y}
\end{equation*}

\bigskip

\begin{itemize}
\item  SOUTH and NORTH boundary:
\end{itemize}

\bigskip

\begin{equation}
S_{m0\alpha}=G_{\alpha y}U_{\alpha}s_{y}\text{,}   \label{c123}
\end{equation}

\begin{equation}
S_{m1\alpha}=-G_{\alpha y}s_{y}\text{,}   \label{c124}
\end{equation}

\begin{equation}
S_{m\alpha}=S_{m0\alpha}+S_{m1\alpha}u_{\alpha}=G_{\alpha y}s_{y}\left(
U_{\alpha}-u_{\alpha}\right) \text{.}   \label{c125}
\end{equation}

\bigskip

\begin{itemize}
\item  WEST and EAST boundary:
\end{itemize}

\bigskip

\begin{equation}
S_{m0\alpha}=G_{\alpha x}U_{\alpha}s_{x}\text{,}   \label{c126}
\end{equation}

\begin{equation}
S_{m1\alpha}=-G_{\alpha x}s_{x}\text{,}   \label{c127}
\end{equation}

\begin{equation}
S_{m\alpha}=S_{m0\alpha}+S_{m1\alpha}u_{\alpha}=G_{\alpha x}s_{x}\left(
U_{\alpha}-u_{\alpha}\right) \text{.}   \label{c128}
\end{equation}

\bigskip

\paragraph{\texttt{BCMOM=5} ... Prescribe the Value of the Momentum Flux per
Unit Area}

\bigskip

\begin{equation*}
\mathtt{MOMPAR(IS,IB,1)}\equiv P_{m\alpha x},\text{ }P_{m\alpha y}
\end{equation*}

\bigskip

\begin{itemize}
\item  SOUTH and NORTH boundary:
\end{itemize}

\bigskip

\begin{equation}
S_{m0\alpha}=P_{m\alpha y}s_{y}\text{,}   \label{c129}
\end{equation}

\begin{equation}
S_{m1\alpha}=0\text{,}   \label{c130}
\end{equation}

\begin{equation}
S_{m\alpha}=S_{m0\alpha}+S_{m1\alpha}u_{\alpha}=P_{m\alpha y}s_{y}\text{.} 
\label{c131}
\end{equation}

\bigskip

\begin{itemize}
\item  WEST and EAST boundary:
\end{itemize}

\bigskip

\begin{equation}
S_{m0\alpha}=P_{m\alpha x}s_{x}\text{,}   \label{c132}
\end{equation}

\begin{equation}
S_{m1\alpha}=0\text{,}   \label{c133}
\end{equation}

\begin{equation}
S_{m\alpha}=S_{m0\alpha}+S_{m1\alpha}u_{\alpha}=P_{m\alpha x}s_{x}\text{.} 
\label{c134}
\end{equation}

\bigskip

\paragraph{\texttt{BCMOM=6} ... Prescribe the Value of the Total Momentum
Flux for a Constant Parallel Velocity (Not Yet Available)}

\bigskip

\bigskip

\paragraph{\texttt{BCMOM=7} ... Prescribe the Density as a Function of the
Other Plasma Parameters (Not Yet Available)}

\bigskip

\bigskip

\paragraph{\texttt{BCMOM=8} ... Limited Shear (Only for the SOUTH and NORTH
Boundary)}

\bigskip

\begin{itemize}
\item  SOUTH boundary:
\end{itemize}

\bigskip

\begin{equation}
a=\min\left( 1\text{, }\sqrt{\frac{\widehat{T}_{i\left( ix,S\right) }+%
\widehat{T}_{i\left( ix+1,S\right) }}{\widehat{T}_{i\left( ix,S+1\right) }+%
\widehat{T}_{i\left( ix+1,S+1\right) }}}\right) \text{,}   \label{c135}
\end{equation}

\begin{equation}
\overline{n}_{\alpha}=\frac{1}{2}\left( n_{\alpha\left( ix,S+1\right)
}+n_{\alpha\left( ix+1,S+1\right) }\right) \text{,}   \label{c136}
\end{equation}

\begin{equation}
\widehat{s}=\frac{1}{2}\left( b_{x\left( ix,S+1\right) }+b_{x\left(
ix+1,S+1\right) }\right) \text{,}   \label{c137}
\end{equation}

\begin{equation}
\widehat{c}=\frac{1}{2}\left( \sqrt{1-b_{x\left( ix,S+1\right) }^{2}}+\sqrt{%
1-b_{x\left( ix+1,S+1\right) }^{2}}\right) \text{,}   \label{c138}
\end{equation}

\begin{equation}
S_{m0\alpha}=\mathcal{V}s_{y}A_{\alpha}\overline{n}_{\alpha}a\left(
u_{\alpha}-\widehat{c}u_{D\alpha}\right) /\widehat{s}\text{,}   \label{c139}
\end{equation}

\begin{equation}
S_{m1\alpha}=-\mathcal{V}s_{y}A_{\alpha}\overline{n}_{\alpha}\text{,} 
\label{c140}
\end{equation}

\begin{equation}
S_{m\alpha}=S_{m0\alpha}+S_{m1\alpha}u_{\alpha}=\mathcal{V}s_{y}A_{\alpha }%
\overline{n}_{\alpha}\left[ a\left( u_{\alpha}-\widehat{c}u_{D\alpha
}\right) /\widehat{s}-u_{\alpha}\right] \text{.}   \label{c141}
\end{equation}

\bigskip

\begin{itemize}
\item  NORTH boundary:
\end{itemize}

\bigskip

\begin{equation}
a=\min\left( 1\text{, }\sqrt{\frac{\widehat{T}_{i\left( ix,N-1\right) }+%
\widehat{T}_{i\left( ix+1,N-1\right) }}{\widehat{T}_{i\left( ix,N\right) }+%
\widehat{T}_{i\left( ix+1,N\right) }}}\right) \text{,}   \label{c142}
\end{equation}

\begin{equation}
\overline{n}_{\alpha}=\frac{1}{2}\left( n_{\alpha\left( ix,N\right)
}+n_{\alpha\left( ix+1,N\right) }\right) \text{,}   \label{c143}
\end{equation}

\begin{equation}
\widehat{s}=\frac{1}{2}\left( b_{x\left( ix,N\right) }+b_{x\left(
ix+1,N\right) }\right) \text{,}   \label{c144}
\end{equation}

\begin{equation}
\widehat{c}=\frac{1}{2}\left( \sqrt{1-b_{x\left( ix,N\right) }^{2}}+\sqrt{%
1-b_{x\left( ix+1,N\right) }^{2}}\right) \text{,}   \label{c145}
\end{equation}

\begin{equation}
S_{m0\alpha}=\mathcal{V}s_{y}A_{\alpha}\overline{n}_{\alpha}a\left(
u_{\alpha}-\widehat{c}u_{D\alpha}\right) /\widehat{s}\text{,}   \label{c146}
\end{equation}

\begin{equation}
S_{m1\alpha}=-\mathcal{V}s_{y}A_{\alpha}\overline{n}_{\alpha}\text{,} 
\label{c147}
\end{equation}

\begin{equation}
S_{m\alpha}=S_{m0\alpha}+S_{m1\alpha}u_{\alpha}=\mathcal{V}s_{y}A_{\alpha }%
\overline{n}_{\alpha}\left[ a\left( u_{\alpha}-\widehat{c}u_{D\alpha
}\right) /\widehat{s}-u_{\alpha}\right] \text{.}   \label{c148}
\end{equation}

\bigskip

\paragraph{\texttt{BCMOM=9} ... Prescribe the Value of the Total Momentum
Flux with Constant Flux Density}

\bigskip

\begin{equation*}
\mathtt{MOMPAR(IS,IB,1)}\equiv\mathcal{S}_{m\alpha x},\text{ }\mathcal{S}%
_{m\alpha y}
\end{equation*}

\bigskip

\begin{itemize}
\item  SOUTH and NORTH boundary:
\end{itemize}

\bigskip

\begin{equation}
s_{tot}=\sum_{j}^{S\text{ or }N}s_{yj}\text{,}   \label{c149}
\end{equation}

\begin{equation}
S_{m0\alpha}=\left( \mathcal{S}_{m\alpha y}/s_{tot}\right) s_{y}\text{,} 
\label{c150}
\end{equation}

\begin{equation}
S_{m1\alpha}=0\text{,}   \label{c151}
\end{equation}

\begin{equation}
S_{m\alpha}=S_{m0\alpha}+S_{m1\alpha}u_{\alpha}=\left( \mathcal{S}_{m\alpha
y}/s_{tot}\right) s_{y}\text{.}   \label{c152}
\end{equation}

\bigskip

\begin{itemize}
\item  WEST and EAST boundary:
\end{itemize}

\bigskip

\begin{equation}
s_{tot}=\sum_{j}^{W\text{ or }E}s_{xj}\text{,}   \label{c153}
\end{equation}

\begin{equation}
S_{m0\alpha}=\left( \mathcal{S}_{m\alpha x}/s_{tot}\right) s_{x}\text{,} 
\label{c154}
\end{equation}

\begin{equation}
S_{m1\alpha}=0\text{,}   \label{c155}
\end{equation}

\begin{equation}
S_{m\alpha}=S_{m0\alpha}+S_{m1\alpha}u_{\alpha}=\left( \mathcal{S}_{m\alpha
x}/s_{tot}\right) s_{x}\text{.}   \label{c156}
\end{equation}

\bigskip

\paragraph{\texttt{BCMOM=10} ... Prescribe the Value of the Decay Length for
the Momentum (Only for the SOUTH and NORTH Boundary)}

\bigskip

\begin{equation*}
\mathtt{MOMPAR(IS,IB,1)}\equiv
l_{u_{\alpha}}=u_{\alpha}/\partial_{y}u_{\alpha}
\end{equation*}

\bigskip

\begin{equation}
S_{m0\alpha}=A_{\alpha}m_{p}u_{\alpha}S_{p0\alpha}\text{,}   \label{c157}
\end{equation}

\begin{equation}
S_{m1\alpha}=-\frac{\eta_{a}^{\alpha}}{l_{u_{\alpha}}}s_{y}+A_{%
\alpha}m_{p}S_{p1\alpha}n_{\alpha}\text{,}   \label{c158}
\end{equation}

\begin{equation}
S_{m\alpha}=S_{m0\alpha}+S_{m1\alpha}u_{\alpha}=A_{\alpha}m_{p}u_{\alpha
}S_{p\alpha}-\frac{\eta_{a}^{\alpha}}{l_{u_{\alpha}}}u_{\alpha}s_{y}\text{.} 
\label{c159}
\end{equation}

\bigskip

\paragraph{\texttt{BCMOM=11} ... Prescribe the Value of the Decay Length for
the Momentum (Only for the SOUTH Boundary, Probably Error in the Source
Code!)}

\bigskip

\begin{equation}
S_{m0\alpha}=\eta_{a}^{\alpha}n_{\alpha}A_{\alpha}\sqrt{1-b_{x}^{2}}%
u_{\alpha }\text{,}   \label{c160}
\end{equation}

\begin{equation}
S_{m1\alpha}=-D_{an}^{\alpha}A_{\alpha}\frac{n_{\alpha\left( S\right)
}-n_{\alpha\left( S+1\right) }}{h_{y1}}-\frac{\eta_{a}^{\alpha}}{h_{y1}}%
n_{\alpha}A_{\alpha}\sqrt{1-b_{x}^{2}}\text{,}   \label{c161}
\end{equation}

\begin{equation}
S_{m\alpha}=S_{m0\alpha}+S_{m1\alpha}u_{\alpha}\text{.}   \label{c162}
\end{equation}

\bigskip

\subsubsection{Boundary Conditions for Electron Energy Equation}

\bigskip

\paragraph{\texttt{BCENE=1} ... Prescribe the Value of the Electron
Temperature}

\bigskip

\begin{equation*}
\mathtt{ENEPAR(IB,1)}\equiv\mathcal{T}_{e}
\end{equation*}

\bigskip

\begin{itemize}
\item  SOUTH and NORTH boundary:
\end{itemize}

\bigskip

\begin{equation}
S_{he0}=\mathcal{VN}_{e}\mathcal{T}_{e}eVs_{y}\text{,}   \label{c163}
\end{equation}

\begin{equation}
S_{he1}=-\mathcal{VN}_{e}s_{y}\text{,}   \label{c164}
\end{equation}

\begin{equation}
S_{he}=S_{he0}+S_{he1}\widehat{T}_{e}=\mathcal{VN}_{e}s_{y}\left( \mathcal{T}%
_{e}eV-\widehat{T}_{e}\right) \text{.}   \label{c165}
\end{equation}

\bigskip

\begin{itemize}
\item  WEST and EAST boundary:
\end{itemize}

\bigskip

\begin{equation}
S_{he0}=\mathcal{UN}_{e}\mathcal{T}_{e}eVs_{x}\text{,}   \label{c166}
\end{equation}

\begin{equation}
S_{he1}=-\mathcal{UN}_{e}s_{x}\text{,}   \label{c167}
\end{equation}

\begin{equation}
S_{he}=S_{he0}+S_{he1}\widehat{T}_{e}=\mathcal{UN}_{e}s_{x}\left( \mathcal{T}%
_{e}eV-\widehat{T}_{e}\right) \text{.}   \label{c168}
\end{equation}

\bigskip

\paragraph{\texttt{BCENE=2} ... Prescribe the Value of the Gradient of the
Electron Temperature}

\bigskip

\begin{equation*}
\mathtt{ENEPAR(IB,1)}\equiv\partial_{x}\mathcal{T}_{e},\text{ }\partial _{y}%
\mathcal{T}_{e}
\end{equation*}

\bigskip

\begin{itemize}
\item  SOUTH boundary:
\end{itemize}

\bigskip

\begin{equation}
l_{y}=\frac{1}{2}\left( h_{y1\left( ix,S\right) }+h_{y1\left( ix,S+1\right)
}\right) \text{,}   \label{c169a}
\end{equation}

\begin{equation}
S_{he0}=\mathcal{VN}_{e}eVs_{y}\left[ \left( -\partial_{y}\mathcal{T}%
_{e}\right) l_{y}+\widehat{T}_{e}\right] \text{,}   \label{c169}
\end{equation}

\begin{equation}
S_{he1}=-\mathcal{VN}_{e}eVs_{y}\text{,}   \label{c170}
\end{equation}

\begin{equation}
S_{he}=S_{he0}+S_{he1}\widehat{T}_{e}=\mathcal{VN}_{e}eVs_{y}\left[ \left(
-\partial_{y}\mathcal{T}_{e}\right) l_{y}\right] \text{.}   \label{c171}
\end{equation}

\bigskip

\begin{itemize}
\item  NORTH boundary:
\end{itemize}

\bigskip

\begin{equation}
l_{y}=\frac{1}{2}\left( h_{y1\left( ix,N-1\right) }+h_{y1\left( ix,N\right)
}\right) \text{,}   \label{c172a}
\end{equation}

\begin{equation}
S_{he0}=\mathcal{VN}_{e}eVs_{y}\left[ \left( \partial_{y}\mathcal{T}%
_{e}\right) l_{y}+\widehat{T}_{e}\right] \text{,}   \label{c172}
\end{equation}

\begin{equation}
S_{he1}=-\mathcal{VN}_{e}eVs_{y}\text{,}   \label{c173}
\end{equation}

\begin{equation}
S_{he}=S_{he0}+S_{he1}\widehat{T}_{e}=\mathcal{VN}_{e}eVs_{y}\left[ \left(
\partial_{y}\mathcal{T}_{e}\right) l_{y}\right] \text{.}   \label{c174}
\end{equation}

\bigskip

\begin{itemize}
\item  WEST boundary:
\end{itemize}

\bigskip

\begin{equation}
l_{x}=\frac{1}{2}\left( h_{x\left( W,iy\right) }+h_{x\left( W+1,iy\right)
}\right) \text{,}   \label{c175a}
\end{equation}

\begin{equation}
S_{he0}=\mathcal{UN}_{e}eVs_{x}\left[ \left( -\partial_{x}\mathcal{T}%
_{e}\right) l_{x}+\widehat{T}_{e}\right] \text{,}   \label{c175}
\end{equation}

\begin{equation}
S_{he1}=-\mathcal{UN}_{e}eVs_{x}\text{,}   \label{c176}
\end{equation}

\begin{equation}
S_{he}=S_{he0}+S_{he1}\widehat{T}_{e}=\mathcal{UN}_{e}eVs_{x}\left[ \left(
-\partial_{x}\mathcal{T}_{e}\right) l_{x}\right] \text{.}   \label{c177}
\end{equation}

\bigskip

\begin{itemize}
\item  EAST boundary:
\end{itemize}

\bigskip

\begin{equation}
l_{x}=\frac{1}{2}\left( h_{x\left( E-1,iy\right) }+h_{x\left( E,iy\right)
}\right) \text{,}   \label{c178a}
\end{equation}

\begin{equation}
S_{he0}=\mathcal{UN}_{e}eVs_{x}\left[ \left( \partial_{x}\mathcal{T}%
_{e}\right) l_{x}+\widehat{T}_{e}\right] \text{,}   \label{c178}
\end{equation}

\begin{equation}
S_{he1}=-\mathcal{UN}_{e}eVs_{x}\text{,}   \label{c179}
\end{equation}

\begin{equation}
S_{he}=S_{he0}+S_{he1}\widehat{T}_{e}=\mathcal{UN}_{e}eVs_{x}\left[ \left(
\partial_{x}\mathcal{T}_{e}\right) l_{x}\right] \text{.}   \label{c180}
\end{equation}

\bigskip

\paragraph{\texttt{BCENE=3} ... Sheath Conditions, Electron Energy
Transmission Coefficient is the Free Input Parameter (Only for the WEST and
EAST Boundary)}

\bigskip

\begin{equation*}
\mathtt{ENEPAR(IB,1)}\equiv\delta_{e}, 
\end{equation*}

\begin{equation*}
\mathtt{MOMPAR(IS,IB,1)}\equiv\mathcal{M}_{\alpha1}, 
\end{equation*}

\begin{equation*}
\mathtt{MOMPAR(1,IB,2)}\equiv\mathcal{M}_{12}, 
\end{equation*}

\begin{equation*}
\mathtt{POTPAR(IB,2)}\equiv\Psi 
\end{equation*}

\bigskip

\begin{itemize}
\item  WEST boundary:
\end{itemize}

\bigskip

IF $\left( \mathcal{M}_{12}<0.5\right) $ THEN:

\begin{equation}
u_{out}=b_{x}\mathcal{M}_{11}c_{s}\text{,}   \label{c181}
\end{equation}

\begin{equation}
\Gamma_{eout}=\min\left( 0\text{, }-u_{out}n_{e}-j_{x}/e_{0}\right) \text{,} 
\label{c182}
\end{equation}

\bigskip

ELSE $\left( \mathcal{M}_{12}\geq0.5\right) $:

\begin{equation}
u_{test\alpha}=u_{\alpha\left( W+1,iy\right) }\left( 1+\frac{h_{x\left(
W+1,iy\right) }}{h_{x\left( W+2,iy\right) }}\right) -u_{\alpha\left(
W+2,iy\right) }\frac{h_{x\left( W+1,iy\right) }}{h_{x\left( W+2,iy\right) }}%
\text{,}   \label{c183}
\end{equation}

\begin{equation}
u_{out\alpha}=\max\left( b_{x}\mathcal{M}_{\alpha1}c_{s\alpha}\text{, }%
-u_{test\alpha}\right) \text{,}   \label{c184}
\end{equation}

\begin{equation}
\Gamma_{ex}=-\sum_{\alpha}u_{out\alpha}Z_{\alpha}n_{\alpha}\text{,} 
\label{c185}
\end{equation}

\begin{equation}
\Gamma_{eout}=\min\left( 0\text{, }\Gamma_{ex}-j_{x}/e_{0}\right) \text{.} 
\label{c186}
\end{equation}

\bigskip

\begin{equation}
S_{he0}=0\text{,}   \label{c187}
\end{equation}

\begin{equation}
S_{he1=}\left[ \delta_{e}+\frac{e_{0}\left( \Phi_{\left( W\right)
}-\Psi\right) }{\widehat{T}_{e}}\right] \Gamma_{eout}s_{x}\text{,} 
\label{c188}
\end{equation}

\begin{equation}
S_{he}=S_{he0}+S_{he1}\widehat{T}_{e}=\left[ \delta_{e}+\frac{e_{0}\left(
\Phi_{\left( W\right) }-\Psi\right) }{\widehat{T}_{e}}\right] \Gamma_{eout}%
\widehat{T}_{e}s_{x}\text{,}   \label{c189}
\end{equation}
where the free input parameter $\Psi$ is the electric potential at the
magnetic presheath edge, i.e., at the boundary between the collisional
presheath and collisionless magnetic presheath.

\bigskip

\begin{itemize}
\item  EAST boundary:
\end{itemize}

\bigskip

IF $\left( \mathcal{M}_{12}<0.5\right) $ THEN:

\begin{equation}
u_{out}=b_{x}\mathcal{M}_{11}c_{s}\text{,}   \label{c190}
\end{equation}

\begin{equation}
\Gamma_{eout}=\max\left( 0\text{, }u_{out}n_{e}-j_{x}/e_{0}\right) \text{,} 
\label{c191}
\end{equation}

\bigskip

ELSE $\left( \mathcal{M}_{12}\geq0.5\right) $:

\begin{equation}
u_{test\alpha}=u_{\alpha\left( E-2,iy\right) }\left( 1+\frac{h_{x\left(
E-1,iy\right) }}{h_{x\left( E-2,iy\right) }}\right) -u_{\alpha\left(
E-3,iy\right) }\frac{h_{x\left( E-1,iy\right) }}{h_{x\left( E-2,iy\right) }}%
\text{,}   \label{c192}
\end{equation}

\begin{equation}
u_{out\alpha}=\max\left( b_{x}\mathcal{M}_{\alpha1}c_{s\alpha}\text{, }%
u_{test\alpha}\right) \text{,}   \label{c193}
\end{equation}

\begin{equation}
\Gamma_{ex}=\sum_{\alpha}u_{out\alpha}Z_{\alpha}n_{\alpha}\text{,} 
\label{c194}
\end{equation}

\begin{equation}
\Gamma_{eout}=\max\left( 0\text{, }\Gamma_{ex}-j_{x}/e_{0}\right) \text{.} 
\label{c195}
\end{equation}

\bigskip

\begin{equation}
S_{he0}=0\text{,}   \label{c196}
\end{equation}

\begin{equation}
S_{he1=}-\left[ \delta_{e}+\frac{e_{0}\left( \Phi_{\left( E\right)
}-\Psi\right) }{\widehat{T}_{e}}\right] \Gamma_{eout}s_{x}\text{,} 
\label{c197}
\end{equation}

\begin{equation}
S_{he}=S_{he0}+S_{he1}\widehat{T}_{e}=-\left[ \delta_{e}+\frac{e_{0}\left(
\Phi_{\left( E\right) }-\Psi\right) }{\widehat{T}_{e}}\right] \Gamma_{eout}%
\widehat{T}_{e}s_{x}\text{,}   \label{c198}
\end{equation}
where the free input parameter $\Psi$ is the electric potential at the
magnetic presheath edge, i.e., at the boundary between the collisional
presheath and collisionless magnetic presheath.

\bigskip

\paragraph{\texttt{BCENE=4} ... Prescribe the Value of the Electron
Temperature ''Weakly'', Using as the Control Parameter Some Prescribed
Electron Flux}

\bigskip

\begin{equation*}
\mathtt{ENEPAR(IB,1)}\equiv\mathcal{T}_{e}, 
\end{equation*}

\begin{equation*}
\mathtt{ENEPAR(IB,2)}\equiv\Gamma_{ex},\text{ }\Gamma_{ey}
\end{equation*}

\bigskip

\begin{itemize}
\item  SOUTH and NORTH boundary:
\end{itemize}

\bigskip

\begin{equation}
S_{he0}=\mathcal{\Gamma}_{ey}\mathcal{T}_{e}s_{y}\text{,}   \label{c199}
\end{equation}

\begin{equation}
S_{he1}=-\left( \mathcal{\Gamma}_{ey}/eV\right) s_{y}\text{,}   \label{c200}
\end{equation}

\begin{equation}
S_{he}=S_{he0}+S_{he1}\widehat{T}_{e}=\mathcal{\Gamma}_{ey}s_{y}\left( 
\mathcal{T}_{e}-\widehat{T}_{e}/eV\right) \text{.}   \label{c201}
\end{equation}

\bigskip

\begin{itemize}
\item  WEST and EAST boundary:
\end{itemize}

\bigskip

\begin{equation}
S_{he0}=\mathcal{\Gamma}_{ex}\mathcal{T}_{e}s_{x}\text{,}   \label{c202}
\end{equation}

\begin{equation}
S_{he1}=-\left( \mathcal{\Gamma}_{ex}/eV\right) s_{x}\text{,}   \label{c203}
\end{equation}

\begin{equation}
S_{he}=S_{he0}+S_{he1}\widehat{T}_{e}=\mathcal{\Gamma}_{ex}s_{x}\left( 
\mathcal{T}_{e}-\widehat{T}_{e}/eV\right) \text{.}   \label{c204}
\end{equation}

\bigskip

\paragraph{\texttt{BCENE=5} ... Prescribe the Value of the Electron Energy
Flux per Unit Area}

\bigskip

\begin{equation*}
\mathtt{ENEPAR(IB,1)}\equiv Q_{hex},Q_{hey}
\end{equation*}

\bigskip

\begin{itemize}
\item  SOUTH and NORTH boundary:
\end{itemize}

\begin{equation}
S_{he0}=Q_{hey}s_{y}\text{,}   \label{c205}
\end{equation}

\begin{equation}
S_{he1}=0\text{,}   \label{c206}
\end{equation}

\begin{equation}
S_{he}=S_{he0}+S_{he1}\widehat{T}_{e}=Q_{hey}s_{y}\text{.}   \label{c207}
\end{equation}

\bigskip

\begin{itemize}
\item  WEST and NORTH boundary:
\end{itemize}

\begin{equation}
S_{he0}=Q_{hex}s_{x}\text{,}   \label{c208}
\end{equation}

\begin{equation}
S_{he1}=0\text{,}   \label{c209}
\end{equation}

\begin{equation}
S_{he}=S_{he0}+S_{he1}\widehat{T}_{e}=Q_{hex}s_{x}\text{.}   \label{c210}
\end{equation}

\bigskip

\paragraph{\texttt{BCENE=6} ... Prescribe the Value of the Total Electron
Energy Flux for a Constant Electron Temperature (Only for the NORTH and
SOUTH Boundary)}

\bigskip

\begin{equation*}
\mathtt{ENEPAR(IB,1)}\equiv\mathcal{Q}
\end{equation*}

\bigskip

The electron temperature $\widehat{T}_{1}$, which is calculated from the
given total electron energy flux $\mathcal{Q}$ (the physics behind that
calculation is unclear to the author of this report), is used in the
calculation of the boundary electron heat source.

\bigskip

\begin{equation}
S_{he0}=\mathcal{VN}_{e}\widehat{T}_{1}s_{y}\text{,}   \label{c211}
\end{equation}

\begin{equation}
S_{he1}=-\mathcal{VN}_{e}s_{y}\text{,}   \label{c212}
\end{equation}

\begin{equation}
S_{he}=S_{he0}+S_{he1}\widehat{T}_{e}=\mathcal{VN}_{e}s_{y}\left( \widehat
{T}_{1}-\widehat{T}_{e}\right) \text{.}   \label{c213}
\end{equation}

\bigskip

\paragraph{\texttt{BCENE=7} ... Prescribe the Electron Temperature as a
Function of Other Plasma Parameters (Not Yet Available)}

\bigskip

\bigskip

\paragraph{\texttt{BCENE=8} ... Prescribe the Value of the Total Electron
Energy Flux with a Constant Flux Density}

\bigskip

\begin{equation*}
\mathtt{ENEPAR(IB,1)}\equiv\QTR{sl}{Q}_{hex}\QTR{sl}{,}\text{ }\QTR{sl}{Q}%
_{hey}
\end{equation*}

\bigskip

\begin{itemize}
\item  SOUTH\ and NORTH boundary:
\end{itemize}

\bigskip

\begin{equation}
s_{tot}=\sum_{j}^{S\text{ or }N}s_{yj}\text{,}   \label{c214}
\end{equation}

\begin{equation}
S_{he0}=\left( \QTR{sl}{Q}_{hey}/s_{tot}\right) s_{y}\text{,}   \label{c215}
\end{equation}

\begin{equation}
S_{he1}=0\text{,}   \label{c216}
\end{equation}

\begin{equation}
S_{he}=S_{he0}+S_{he1}\widehat{T}_{e}=\left( \QTR{sl}{Q}_{hey}/s_{tot}%
\right) s_{y}\text{.}   \label{c217}
\end{equation}

\bigskip

\begin{itemize}
\item  WEST and NORTH boundary:
\end{itemize}

\bigskip

\begin{equation}
s_{tot}=\sum_{j}^{W\text{ or }E}s_{xj}\text{,}   \label{c218}
\end{equation}

\begin{equation}
S_{he0}=\left( \QTR{sl}{Q}_{hex}/s_{tot}\right) s_{x}\text{,}   \label{c219}
\end{equation}

\begin{equation}
S_{he1}=0\text{,}   \label{c220}
\end{equation}

\begin{equation}
S_{he}=S_{he0}+S_{he1}\widehat{T}_{e}=\left( \QTR{sl}{Q}_{hex}/s_{tot}%
\right) s_{x}\text{.}   \label{c221}
\end{equation}

\bigskip

\paragraph{\texttt{BCENE=9} ... Prescribe the Value of the Decay Length for
the Electron Temperature (Only for the SOUTH and NORTH Boundary)}

\bigskip

\begin{equation*}
\mathtt{ENEPAR(IB,1)}\equiv l_{T_{e}}=\widehat{T}_{e}/\partial_{y}\widehat
{T}_{e}
\end{equation*}

\bigskip

\begin{equation}
S_{he0}=\sum_{\alpha}Z_{\alpha}\frac{5}{2}\widehat{T}_{e}S_{p0\alpha}\text{,}
\label{c222}
\end{equation}

\begin{equation}
S_{he1}=-\frac{\kappa_{a}^{e}}{l_{T_{e}}}s_{y}+\sum_{\alpha}Z_{\alpha}\frac
{5}{2}S_{p1\alpha}n_{\alpha}\text{,}   \label{c223}
\end{equation}

\begin{equation}
S_{he}=S_{he0}+S_{he1}\widehat{T}_{e}=-\frac{\kappa_{a}^{e}}{l_{T_{e}}}%
\widehat{T}_{e}s_{y}+\sum_{\alpha}Z_{\alpha}\frac{5}{2}\widehat{T}%
_{e}S_{p\alpha}\text{.}   \label{c224}
\end{equation}

\bigskip

\paragraph{\texttt{BCENE=10} ... Feedback Option for the Core (Only for the
SOUTH Boundary)}

\bigskip

\begin{equation}
S_{he0}=C_{he0}s_{y}\text{,}   \label{c225}
\end{equation}

\begin{equation}
S_{he1}=0\text{,}   \label{c226}
\end{equation}

\begin{equation}
S_{he}=S_{he0}+S_{he1}\widehat{T}_{e}=C_{he0}s_{y}\text{.}   \label{c227}
\end{equation}

\bigskip

\subsubsection{Boundary Conditions for Ion Energy Equation}

\bigskip

\paragraph{\texttt{BCENI=1} ... Prescribe the Value of the Ion Temperature}

\bigskip

\begin{equation*}
\mathtt{ENIPAR(IB,1)}\equiv\mathcal{T}_{i}
\end{equation*}

\bigskip

\begin{itemize}
\item  SOUTH and NORTH boundary:
\end{itemize}

\bigskip

\begin{equation}
S_{hi0}=\mathcal{VN}_{e}\mathcal{T}_{i}eVs_{y}\text{,}   \label{c228}
\end{equation}

\begin{equation}
S_{hi1}=-\mathcal{VN}_{e}s_{y}\text{,}   \label{c229}
\end{equation}

\begin{equation}
S_{hi}=S_{hi0}+S_{hi1}\widehat{T}_{e}=\mathcal{VN}_{e}s_{y}\left( \mathcal{T}%
_{i}eV-\widehat{T}_{i}\right) \text{.}   \label{c230}
\end{equation}

\bigskip

\begin{itemize}
\item  WEST and EAST boundary:
\end{itemize}

\bigskip

\begin{equation}
S_{hi0}=\mathcal{UN}_{e}\mathcal{T}_{i}eVs_{x}\text{,}   \label{c231}
\end{equation}

\begin{equation}
S_{hi1}=-\mathcal{UN}_{e}s_{x}\text{,}   \label{c232}
\end{equation}

\begin{equation}
S_{hi}=S_{hi0}+S_{hi1}\widehat{T}_{i}=\mathcal{UN}_{e}s_{x}\left( \mathcal{T}%
_{i}eV-\widehat{T}_{i}\right) \text{.}   \label{c233}
\end{equation}

\bigskip

\paragraph{\texttt{BCENI=2} ... Prescribe the Value of the Gradient of the
Ion Temperature}

\bigskip

\begin{equation*}
\mathtt{ENIPAR(IB,1)}\equiv\partial_{x}\mathcal{T}_{i},\text{ }\partial _{y}%
\mathcal{T}_{i}
\end{equation*}

\bigskip

\begin{itemize}
\item  SOUTH boundary:
\end{itemize}

\bigskip

\begin{equation}
l_{y}=\frac{1}{2}\left( h_{y1\left( ix,S\right) }+h_{y1\left( ix,S+1\right)
}\right) \text{,}   \label{c234a}
\end{equation}

\begin{equation}
S_{hi0}=\mathcal{VN}_{e}eVs_{y}\left[ \left( -\partial_{y}\mathcal{T}%
_{i}\right) l_{y}+\widehat{T}_{i}\right] \text{,}   \label{c234}
\end{equation}

\begin{equation}
S_{hi1}=-\mathcal{VN}_{e}eVs_{y}\text{,}   \label{c235}
\end{equation}

\begin{equation}
S_{hi}=S_{hi0}+S_{hi1}\widehat{T}_{i}=\mathcal{VN}_{e}eVs_{y}\left[ \left(
-\partial_{y}\mathcal{T}_{i}\right) l_{y}\right] \text{.}   \label{c236}
\end{equation}

\bigskip

\begin{itemize}
\item  NORTH boundary:
\end{itemize}

\bigskip

\begin{equation}
l_{y}=\frac{1}{2}\left( h_{y1\left( ix,N-1\right) }+h_{y1\left( ix,N\right)
}\right) \text{,}   \label{c237a}
\end{equation}

\begin{equation}
S_{hi0}=\mathcal{VN}_{e}eVs_{y}\left[ \left( \partial_{y}\mathcal{T}%
_{i}\right) l_{y}+\widehat{T}_{i}\right] \text{,}   \label{c237}
\end{equation}

\begin{equation}
S_{hi1}=-\mathcal{VN}_{e}eVs_{y}\text{,}   \label{c238}
\end{equation}

\begin{equation}
S_{hi}=S_{hi0}+S_{hi1}\widehat{T}_{i}=\mathcal{VN}_{e}eVs_{y}\left[ \left(
\partial_{y}\mathcal{T}_{i}\right) l_{y}\right] \text{.}   \label{c239}
\end{equation}

\bigskip

\begin{itemize}
\item  WEST boundary:
\end{itemize}

\bigskip

\begin{equation}
l_{x}=\frac{1}{2}\left( h_{x\left( W,iy\right) }+h_{x\left( W+1,iy\right)
}\right) \text{,}   \label{c240a}
\end{equation}

\begin{equation}
S_{hi0}=\mathcal{UN}_{e}eVs_{x}\left[ \left( -\partial_{x}\mathcal{T}%
_{i}\right) l_{x}+\widehat{T}_{i}\right] \text{,}   \label{c240}
\end{equation}

\begin{equation}
S_{hi1}=-\mathcal{UN}_{e}eVs_{x}\text{,}   \label{c241}
\end{equation}

\begin{equation}
S_{hi}=S_{hi0}+S_{hi1}\widehat{T}_{i}=\mathcal{UN}_{e}eVs_{x}\left[ \left(
-\partial_{x}\mathcal{T}_{i}\right) l_{x}\right] \text{.}   \label{c242}
\end{equation}

\bigskip

\begin{itemize}
\item  EAST boundary:
\end{itemize}

\bigskip

\begin{equation}
l_{x}=\frac{1}{2}\left( h_{x\left( E-1,iy\right) }+h_{x\left( E,iy\right)
}\right) \text{,}   \label{c243a}
\end{equation}

\begin{equation}
S_{hi0}=\mathcal{UN}_{e}eVs_{x}\left[ \left( \partial_{x}\mathcal{T}%
_{i}\right) l_{x}+\widehat{T}_{i}\right] \text{,}   \label{c243}
\end{equation}

\begin{equation}
S_{hi1}=-\mathcal{UN}_{e}eVs_{x}\text{,}   \label{c244}
\end{equation}

\begin{equation}
S_{hi}=S_{hi0}+S_{hi1}\widehat{T}_{i}=\mathcal{UN}_{e}eVs_{x}\left[ \left(
\partial_{x}\mathcal{T}_{i}\right) l_{x}\right] \text{.}   \label{c245}
\end{equation}

\bigskip

\paragraph{\texttt{BCENI=3} ... Sheath Conditions, Ion Energy Transmission
Coefficients Are the Free Input Parameters (Only for the WEST and EAST
Boundary)}

\bigskip

\begin{equation*}
\mathtt{ENIPAR(IB,1)}\equiv\delta_{i1}, 
\end{equation*}

\begin{equation*}
\mathtt{ENIPAR(IB,2)}\equiv\delta_{i2}, 
\end{equation*}

\begin{equation*}
\mathtt{MOMPAR(IS,IB,1)}\equiv\mathcal{M}_{\alpha1}, 
\end{equation*}

\begin{equation*}
\mathtt{MOMPAR(1,IB,2)}\equiv\mathcal{M}_{12}, 
\end{equation*}

\begin{equation*}
\mathtt{POTPAR(IB,2)}\equiv\Psi 
\end{equation*}

\bigskip

\begin{itemize}
\item  WEST boundary:
\end{itemize}

\bigskip

IF $\left( \mathcal{M}_{12}<0.5\right) $ THEN:

\begin{equation}
u_{out}=b_{x}\mathcal{M}_{11}c_{s}\text{,}   \label{c246}
\end{equation}

\begin{equation}
u_{p}=\frac{u_{out}}{b_{x}}+\sqrt{\frac{1}{b_{x}^{2}}-1}u_{D1}\text{,} 
\label{c247}
\end{equation}

\begin{equation}
S_{hi0}=-\sum_{\alpha}\delta_{i2}\frac{1}{2}u_{out}A_{\alpha}m_{p}n_{\alpha
}\left( u_{p}^{2}+u_{D\alpha}^{2}\right) s_{x}\text{,}   \label{c248}
\end{equation}

\begin{equation}
S_{hi1}=-\sum_{\alpha}u_{out}\delta_{i1}n_{\alpha}s_{x}\text{,} 
\label{c249}
\end{equation}

\begin{equation}
S_{hi}=S_{hi0}+S_{hi1}\widehat{T}_{i}=-\sum_{\alpha}n_{\alpha}u_{out}s_{x}%
\left[ \delta_{i2}\frac{1}{2}A_{\alpha}m_{p}\left(
u_{p}^{2}+u_{D\alpha}^{2}\right) +\delta_{i1}\widehat{T}_{i}\right] \text{,} 
\label{c250}
\end{equation}

\bigskip

ELSE $\left( \mathcal{M}_{12}\geq0.5\right) $:

\begin{equation}
u_{test\alpha}=u_{\alpha\left( W+1,iy\right) }\left( 1+\frac{h_{x\left(
W+1,iy\right) }}{h_{x\left( W+2,iy\right) }}\right) -u_{\alpha\left(
W+2,iy\right) }\frac{h_{x\left( W+1,iy\right) }}{h_{x\left( W+2,iy\right) }}%
\text{,}   \label{c251}
\end{equation}

\begin{equation}
u_{out\alpha}=\max\left( b_{x}\mathcal{M}_{\alpha1}c_{s\alpha}\text{, }%
-u_{test\alpha}\right) \text{,}   \label{cc251}
\end{equation}

\begin{equation}
u_{p\alpha}=\frac{u_{out\alpha}}{b_{x}}+\sqrt{\frac{1}{b_{x}^{2}}-1}%
u_{D\alpha}\text{,}   \label{c252}
\end{equation}

\begin{equation}
S_{hi0}=-\sum_{\alpha}\delta_{i2}\frac{1}{2}u_{out\alpha}A_{\alpha}m_{p}n_{%
\alpha}\left( u_{p\alpha}^{2}+u_{D\alpha}^{2}\right) s_{x}\text{,} 
\label{c253}
\end{equation}

\begin{equation}
S_{hi1}=-\sum_{\alpha}u_{out\alpha}\delta_{i1}n_{\alpha}s_{x}\text{,} 
\label{c254}
\end{equation}

\begin{equation}
S_{hi}=S_{hi0}+S_{hi1}\widehat{T}_{i}=-\sum_{\alpha}n_{\alpha}u_{out\alpha
}s_{x}\left[ \delta_{i2}\frac{1}{2}A_{\alpha}m_{p}\left(
u_{p\alpha}^{2}+u_{D\alpha}^{2}\right) +\delta_{i1}\widehat{T}_{i}\right] 
\text{.}   \label{c255}
\end{equation}

\bigskip

\begin{itemize}
\item  EAST boundary:
\end{itemize}

\bigskip

IF $\left( \mathcal{M}_{12}<0.5\right) $ THEN:

\begin{equation}
u_{out}=b_{x}\mathcal{M}_{11}c_{s}\text{,}   \label{c256}
\end{equation}

\begin{equation}
u_{p}=\frac{u_{out}}{b_{x}}-\sqrt{\frac{1}{b_{x}^{2}}-1}u_{D1}\text{,} 
\label{c257}
\end{equation}

\begin{equation}
S_{hi0}=-\sum_{\alpha}\delta_{i2}\frac{1}{2}u_{out}A_{\alpha}m_{p}n_{\alpha
}\left( u_{p}^{2}+u_{D\alpha}^{2}\right) s_{x}\text{,}   \label{c258}
\end{equation}

\begin{equation}
S_{hi1}=-\sum_{\alpha}u_{out}\delta_{i1}n_{\alpha}s_{x}\text{,} 
\label{c259}
\end{equation}

\begin{equation}
S_{hi}=S_{hi0}+S_{hi1}\widehat{T}_{i}=-\sum_{\alpha}n_{\alpha}u_{out}s_{x}%
\left[ \delta_{i2}\frac{1}{2}A_{\alpha}m_{p}\left(
u_{p}^{2}+u_{D\alpha}^{2}\right) +\delta_{i1}\widehat{T}_{i}\right] \text{,} 
\label{c260}
\end{equation}

\bigskip

ELSE $\left( \mathcal{M}_{12}\geq0.5\right) $:

\begin{equation}
u_{test\alpha}=u_{\alpha\left( E-2,iy\right) }\left( 1+\frac{h_{x\left(
E-1,iy\right) }}{h_{x\left( E-2,iy\right) }}\right) -u_{\alpha\left(
E-3,iy\right) }\frac{h_{x\left( E-1,iy\right) }}{h_{x\left( E-2,iy\right) }}%
\text{,}   \label{c261}
\end{equation}

\begin{equation}
u_{out\alpha}=\max\left( b_{x}\mathcal{M}_{\alpha1}c_{s\alpha}\text{, }%
u_{test\alpha}\right) \text{,}   \label{c262}
\end{equation}

\begin{equation}
u_{p\alpha}=\frac{u_{out\alpha}}{b_{x}}-\sqrt{\frac{1}{b_{x}^{2}}-1}%
u_{D\alpha}\text{,}   \label{c263}
\end{equation}

\begin{equation}
S_{hi0}=-\sum_{\alpha}\delta_{i2}\frac{1}{2}u_{out\alpha}A_{\alpha}m_{p}n_{%
\alpha}\left( u_{p\alpha}^{2}+u_{D\alpha}^{2}\right) s_{x}\text{,} 
\label{c264}
\end{equation}

\begin{equation}
S_{hi1}=-\sum_{\alpha}u_{out\alpha}\delta_{i1}n_{\alpha}s_{x}\text{,} 
\label{c265}
\end{equation}

\begin{equation}
S_{hi}=S_{hi0}+S_{hi1}\widehat{T}_{i}=-\sum_{\alpha}n_{\alpha}u_{out\alpha
}s_{x}\left[ \delta_{i2}\frac{1}{2}A_{\alpha}m_{p}\left(
u_{p\alpha}^{2}+u_{D\alpha}^{2}\right) +\delta_{i1}\widehat{T}_{i}\right] 
\text{.}   \label{c266}
\end{equation}

\bigskip

\paragraph{\texttt{BCENI=4} ... Prescribe the Value of the Ion Temperature
''Weakly'', Using as the Control Parameter Some Prescribed Ion Flux}

\bigskip

\begin{equation*}
\mathtt{ENIPAR(IB,1)}\equiv\mathcal{T}_{i}, 
\end{equation*}

\begin{equation*}
\mathtt{ENIPAR(IB,2)}\equiv\Gamma_{ix},\text{ }\Gamma_{iy}
\end{equation*}

\bigskip

\begin{itemize}
\item  SOUTH and NORTH boundary:
\end{itemize}

\bigskip

\begin{equation}
S_{hi0}=\mathcal{\Gamma}_{iy}\mathcal{T}_{i}s_{y}\text{,}   \label{c267}
\end{equation}

\begin{equation}
S_{hi1}=-\left( \mathcal{\Gamma}_{iy}/eV\right) s_{y}\text{,}   \label{c268}
\end{equation}

\begin{equation}
S_{hi}=S_{hi0}+S_{hi1}\widehat{T}_{i}=\mathcal{\Gamma}_{iy}s_{y}\left( 
\mathcal{T}_{i}-\widehat{T}_{i}/eV\right) \text{.}   \label{c269}
\end{equation}

\bigskip

\begin{itemize}
\item  WEST and EAST boundary:
\end{itemize}

\bigskip

\begin{equation}
S_{hi0}=\mathcal{\Gamma}_{ix}\mathcal{T}_{i}s_{x}\text{,}   \label{c270}
\end{equation}

\begin{equation}
S_{hi1}=-\left( \mathcal{\Gamma}_{ix}/eV\right) s_{x}\text{,}   \label{c271}
\end{equation}

\begin{equation}
S_{hi}=S_{hi0}+S_{hi1}\widehat{T}_{i}=\mathcal{\Gamma}_{ix}s_{x}\left( 
\mathcal{T}_{i}-\widehat{T}_{i}/eV\right) \text{.}   \label{c272}
\end{equation}

\bigskip

\paragraph{\texttt{BCENI=5} ... Prescribe the Value of the Ion Energy Flux
per Unit Area}

\bigskip

\begin{equation*}
\mathtt{ENIPAR(IB,1)}\equiv Q_{hix},Q_{hiy}
\end{equation*}

\bigskip

\begin{itemize}
\item  SOUTH and NORTH boundary:
\end{itemize}

\begin{equation}
S_{hi0}=Q_{hiy}s_{y}\text{,}   \label{c273}
\end{equation}

\begin{equation}
S_{hi1}=0\text{,}   \label{c274}
\end{equation}

\begin{equation}
S_{hi}=S_{hi0}+S_{hi1}\widehat{T}_{i}=Q_{hiy}s_{y}\text{.}   \label{c275}
\end{equation}

\bigskip

\begin{itemize}
\item  WEST and NORTH boundary:
\end{itemize}

\begin{equation}
S_{hi0}=Q_{hix}s_{x}\text{,}   \label{c276}
\end{equation}

\begin{equation}
S_{hi1}=0\text{,}   \label{c277}
\end{equation}

\begin{equation}
S_{hi}=S_{hi0}+S_{hi1}\widehat{T}_{i}=Q_{hix}s_{x}\text{.}   \label{c278}
\end{equation}

\bigskip

\paragraph{\texttt{BCENI=6} ... Prescribe the Value of the Total Ion Energy
Flux for a Constant Ion Temperature (Only for the NORTH and SOUTH Boundary)}

\bigskip

\begin{equation*}
\mathtt{ENIPAR(IB,1)}\equiv\mathcal{Q}
\end{equation*}

\bigskip

The ion temperature $\widehat{T}_{1}$, which is calculated from the given
total ion energy flux $\mathcal{Q}$ (the physics behind that calculation is
unclear to the author of this report), is used in the calculation of the
boundary ion heat source.

\bigskip

\begin{equation}
S_{hi0}=\mathcal{VN}_{e}\widehat{T}_{1}s_{y}\text{,}   \label{c279}
\end{equation}

\begin{equation}
S_{hi1}=-\mathcal{VN}_{e}s_{y}\text{,}   \label{c280}
\end{equation}

\begin{equation}
S_{hi}=S_{hi0}+S_{hi1}\widehat{T}_{i}=\mathcal{VN}_{e}s_{y}\left( \widehat
{T}_{1}-\widehat{T}_{i}\right) \text{.}   \label{c281}
\end{equation}

\bigskip

\paragraph{\texttt{BCENI=7} ... Prescribe the Ion Temperature as a Function
of Other Plasma Parameters (Not Yet Available)}

\bigskip

\bigskip

\paragraph{\texttt{BCENI=8} ... Prescribe the Value of the Total Ion Energy
Flux with a Constant Flux Density}

\bigskip

\begin{equation*}
\mathtt{ENIPAR(IB,1)}\equiv\QTR{sl}{Q}_{hix}\QTR{sl}{,}\text{ }\QTR{sl}{Q}%
_{hiy}
\end{equation*}

\bigskip

\begin{itemize}
\item  SOUTH\ and NORTH boundary:
\end{itemize}

\bigskip

\begin{equation}
s_{tot}=\sum_{j}^{S\text{ or }N}s_{yj}\text{,}   \label{c282}
\end{equation}

\begin{equation}
S_{hi0}=\left( \QTR{sl}{Q}_{hiy}/s_{tot}\right) s_{y}\text{,}   \label{c283}
\end{equation}

\begin{equation}
S_{hi1}=0\text{,}   \label{c284}
\end{equation}

\begin{equation}
S_{hi}=S_{hi0}+S_{hi1}\widehat{T}_{i}=\left( \QTR{sl}{Q}_{hiy}/s_{tot}%
\right) s_{y}\text{.}   \label{c285}
\end{equation}

\bigskip

\begin{itemize}
\item  WEST and NORTH boundary:
\end{itemize}

\bigskip

\begin{equation}
s_{tot}=\sum_{j}^{W\text{ or }E}s_{xj}\text{,}   \label{c286}
\end{equation}

\begin{equation}
S_{hi0}=\left( \QTR{sl}{Q}_{hix}/s_{tot}\right) s_{x}\text{,}   \label{c287}
\end{equation}

\begin{equation}
S_{hi1}=0\text{,}   \label{c288}
\end{equation}

\begin{equation}
S_{hi}=S_{hi0}+S_{hi1}\widehat{T}_{i}=\left( \QTR{sl}{Q}_{hix}/s_{tot}%
\right) s_{x}\text{.}   \label{c289}
\end{equation}

\bigskip

\paragraph{\texttt{BCENI=9} ... Prescribe the Value of the Decay Length for
the Ion Temperature (Only for the SOUTH and NORTH Boundary)}

\bigskip

\begin{equation*}
\mathtt{ENIPAR(IB,1)}\equiv l_{T_{i}}=\widehat{T}_{i}/\partial_{y}\widehat
{T}_{i}
\end{equation*}

\bigskip

\begin{equation}
S_{hi0}=\sum_{\alpha}Z_{\alpha}\frac{5}{2}\widehat{T}_{i}S_{p0\alpha}\text{,}
\label{c290}
\end{equation}

\begin{equation}
S_{hi1}=-\frac{\kappa_{a}^{i}}{l_{T_{i}}}s_{y}+\sum_{\alpha}Z_{\alpha}\frac
{5}{2}S_{p1\alpha}n_{\alpha}\text{,}   \label{c291}
\end{equation}

\begin{equation}
S_{hi}=S_{hi0}+S_{hi1}\widehat{T}_{i}=-\frac{\kappa_{a}^{i}}{l_{T_{i}}}%
\widehat{T}_{i}s_{y}+\sum_{\alpha}Z_{\alpha}\frac{5}{2}\widehat{T}%
_{i}S_{p\alpha}\text{.}   \label{c292}
\end{equation}

\bigskip

\paragraph{\texttt{BCENI=10} ... Feedback Option for the Core (Only for the
SOUTH Boundary)}

\bigskip

\begin{equation}
S_{hi0}=C_{hi0}s_{y}\text{,}   \label{c293}
\end{equation}

\begin{equation}
S_{hi1}=0\text{,}   \label{c294}
\end{equation}

\begin{equation}
S_{hi}=S_{hi0}+S_{hi1}\widehat{T}_{i}=C_{hi0}s_{y}\text{.}   \label{c295}
\end{equation}

\bigskip

\bigskip

\subsubsection{Boundary Conditions for the Electric Potential Equation}

\bigskip

\paragraph{\texttt{BCPOT=1} ... Prescribe the Value of the Potential}

\bigskip

\begin{equation*}
\mathtt{POTPAR(IB,1)}\equiv\Psi 
\end{equation*}

\bigskip

\begin{itemize}
\item  SOUTH and NORTH boundary:
\end{itemize}

\bigskip

\begin{equation}
S_{ch0}=\mathcal{V}\Psi s_{y}\text{,}   \label{c296}
\end{equation}

\begin{equation}
S_{ch1}=-\mathcal{V}s_{y}\text{,}   \label{c297}
\end{equation}

\begin{equation}
S_{ch}=S_{ch0}+S_{ch1}\Phi=\mathcal{V}s_{y}\left( \Psi-\Phi\right) \text{.} 
\label{c298}
\end{equation}

\bigskip

\begin{itemize}
\item  WEST and EAST boundary:
\end{itemize}

\bigskip

\begin{equation}
S_{ch0}=\mathcal{U}\Psi s_{x}\text{,}   \label{c299}
\end{equation}

\begin{equation}
S_{ch1}=-\mathcal{U}s_{x}\text{,}   \label{c300}
\end{equation}

\begin{equation}
S_{ch}=S_{ch0}+S_{ch1}\Phi=\mathcal{U}s_{x}\left( \Psi-\Phi\right) \text{,} 
\label{c301}
\end{equation}
where $\Psi$ is the potential at the magnetic presheath edge and $\Phi$ is
the potential at the boundary.

\bigskip

\paragraph{\texttt{BCPOT=2} ... Prescribe the Value of the Gradient of the
Potential}

\bigskip

\begin{equation*}
\mathtt{POTPAR(IB,1)}\equiv\partial_{x}\Psi,\text{ }\partial_{y}\Psi 
\end{equation*}

\bigskip

\begin{itemize}
\item  SOUTH boundary:
\end{itemize}

\bigskip

\begin{equation}
l_{y}=\frac{1}{2}\left( h_{y1\left( ix,S\right) }+h_{y1\left( ix,S+1\right)
}\right) \text{,}   \label{c302a}
\end{equation}

\begin{equation}
S_{ch0}=\mathcal{V}s_{y}\left[ \left( -\partial_{y}\Psi\right) l_{y}+\Phi%
\right] \text{,}   \label{c302}
\end{equation}

\begin{equation}
S_{ch1}=-\mathcal{V}s_{y}\text{,}   \label{c303}
\end{equation}

\begin{equation}
S_{ch}=S_{ch0}+S_{ch1}\Phi=\mathcal{V}s_{y}\left[ \left(
-\partial_{y}\Psi\right) l_{y}\right] \text{.}   \label{c304}
\end{equation}

\bigskip

\begin{itemize}
\item  NORTH boundary:
\end{itemize}

\bigskip

\begin{equation}
l_{y}=\frac{1}{2}\left( h_{y1\left( ix,N-1\right) }+h_{y1\left( ix,N\right)
}\right) \text{,}   \label{c305a}
\end{equation}

\begin{equation}
S_{ch0}=\mathcal{V}s_{y}\left[ \left( \partial_{y}\Psi\right) l_{y}+\Phi%
\right] \text{,}   \label{c305}
\end{equation}

\begin{equation}
S_{ch1}=-\mathcal{V}s_{y}\text{,}   \label{c306}
\end{equation}

\begin{equation}
S_{ch}=S_{ch0}+S_{ch1}\Phi=\mathcal{V}s_{y}\left[ \left(
\partial_{y}\Psi\right) l_{y}\right] \text{.}   \label{c307}
\end{equation}

\bigskip

\begin{itemize}
\item  WEST boundary:
\end{itemize}

\bigskip

\begin{equation}
l_{x}=\frac{1}{2}\left( h_{x\left( W,iy\right) }+h_{x\left( W+1,iy\right)
}\right) \text{,}   \label{c308a}
\end{equation}

\begin{equation}
S_{ch0}=\mathcal{U}s_{x}\left[ \left( -\partial_{x}\Psi\right) l_{x}+\Phi%
\right] \text{,}   \label{c308}
\end{equation}

\begin{equation}
S_{ch1}=-\mathcal{U}s_{x}\text{,}   \label{c309}
\end{equation}

\begin{equation}
S_{ch}=S_{ch0}+S_{ch1}\Phi=\mathcal{U}s_{x}\left[ \left(
-\partial_{x}\Psi\right) l_{x}\right] \text{.}   \label{c310}
\end{equation}

\bigskip

\begin{itemize}
\item  EAST boundary:
\end{itemize}

\bigskip

\begin{equation}
l_{x}=\frac{1}{2}\left( h_{x\left( E-1,iy\right) }+h_{x\left( E,iy\right)
}\right) \text{,}   \label{c311a}
\end{equation}

\begin{equation}
S_{ch0}=\mathcal{U}s_{x}\left[ \left( \partial_{x}\Psi\right) l_{x}+\Phi%
\right] \text{,}   \label{c311}
\end{equation}

\begin{equation}
S_{ch1}=-\mathcal{U}s_{x}\text{,}   \label{c312}
\end{equation}

\begin{equation}
S_{ch}=S_{ch0}+S_{ch1}\Phi=\mathcal{U}s_{x}\left[ \left(
\partial_{x}\Psi\right) l_{x}\right] \text{.}   \label{c313}
\end{equation}

\bigskip

\paragraph{\texttt{BCPOT=3} ... Sheath Conditions (Only for the WEST and
EAST Boundary)}

\bigskip

\begin{equation*}
\mathtt{POTPAR(IB,2)}\equiv\Psi 
\end{equation*}

\bigskip

\begin{itemize}
\item  WEST boundary:
\end{itemize}

\begin{equation}
j_{i}=-e_{0}\sum_{\alpha}Z_{\alpha}\Gamma_{\alpha x}\text{,}   \label{c314}
\end{equation}

\begin{equation}
j_{e}=\frac{1}{\sqrt{2\pi}}e_{0}n_{e}v_{th,e}S_{1}\exp\left[ -\frac
{e_{0}\left( \Phi-\Psi\right) }{\widehat{T}_{e}}\right] \text{,} 
\label{c315}
\end{equation}

\begin{equation}
S_{ch0}=\left( 1-\gamma_{se}\right) j_{e}-j_{i}+\max\left( j_{i},\left(
1-\gamma_{se}\right) j_{e}\right) e_{0}\Phi/\widehat{T}_{e}\text{,} 
\label{c316}
\end{equation}

\begin{equation}
S_{ch1}=-\max\left( j_{i},\left( 1-\gamma_{se}\right) j_{e}\right) e_{0}/%
\widehat{T}_{e}\text{,}   \label{c317}
\end{equation}

\begin{equation}
S_{ch}=S_{ch0}+S_{ch1}\Phi=\left( 1-\gamma_{se}\right) j_{e}-j_{i}\text{.} 
\label{c318}
\end{equation}

\bigskip

\begin{itemize}
\item  EAST boundary:
\end{itemize}

\begin{equation}
j_{i}=e_{0}\sum_{\alpha}Z_{\alpha}\Gamma_{\alpha x}\text{,}   \label{c319}
\end{equation}

\begin{equation}
j_{e}=\frac{1}{\sqrt{2\pi}}e_{0}n_{e}v_{th,e}S_{1}\exp\left[ -\frac
{e_{0}\left( \Phi-\Psi\right) }{\widehat{T}_{e}}\right] \text{,} 
\label{c320}
\end{equation}

\begin{equation}
S_{ch0}=\left( 1-\gamma_{se}\right) j_{e}-j_{i}+\max\left( j_{i},\left(
1-\gamma_{se}\right) j_{e}\right) e_{0}\Phi/\widehat{T}_{e}\text{,} 
\label{c321}
\end{equation}

\begin{equation}
S_{ch1}=-\max\left( j_{i},\left( 1-\gamma_{se}\right) j_{e}\right) e_{0}/%
\widehat{T}_{e}\text{,}   \label{c322}
\end{equation}

\begin{equation}
S_{ch}=S_{ch0}+S_{ch1}\Phi=\left( 1-\gamma_{se}\right) j_{e}-j_{i}\text{.} 
\label{c323}
\end{equation}

\bigskip

\bigskip

\paragraph{\texttt{BCPOT=4} ... Prescribe the Value of the Potential
''Weakly'', Using the Prescribed Control Parameter, Which Corresponds to the
Specific Electric Conductivity}

\bigskip

\begin{equation*}
\mathtt{POTPAR(IB,1)}\equiv\Psi\text{,}
\end{equation*}

\begin{equation*}
\mathtt{POTPAR(IB,2)}\equiv\sigma_{x}\text{, }\sigma_{y}
\end{equation*}

\bigskip

\begin{itemize}
\item  SOUTH and NORTH boundary:
\end{itemize}

\begin{equation}
S_{ch0}=\sigma_{y}\Psi s_{y}\text{,}   \label{c324}
\end{equation}

\begin{equation}
S_{ch1}=-\sigma_{y}s_{y}\text{,}   \label{c325}
\end{equation}

\begin{equation}
S_{ch}=S_{ch0}+S_{ch1}\Phi=\sigma_{y}s_{y}\left( \Psi-\Phi\right) \text{.} 
\label{c326}
\end{equation}

\bigskip

\begin{itemize}
\item  WEST and EAST boundary:
\end{itemize}

\begin{equation}
S_{ch0}=\sigma_{x}\Psi s_{x}\text{,}   \label{c327}
\end{equation}

\begin{equation}
S_{ch1}=-\sigma_{x}s_{x}\text{,}   \label{c328}
\end{equation}

\begin{equation}
S_{ch}=S_{ch0}+S_{ch1}\Phi=\sigma_{x}s_{x}\left( \Psi-\Phi\right) \text{.} 
\label{c329}
\end{equation}

\bigskip

\bigskip

\paragraph{\texttt{BCPOT=5} ... Prescribe the Value of the Current Flux
Density per Unit Area}

\bigskip

\begin{equation}
\mathtt{POTPAR(IB,1)}\equiv j_{x}\text{, }j_{y}   \label{c330}
\end{equation}

\bigskip

\begin{itemize}
\item  SOUTH and NORTH boundary:
\end{itemize}

\begin{equation}
S_{ch0}=j_{y}s_{y}\text{,}   \label{c331}
\end{equation}

\begin{equation}
S_{ch1}=0\text{,}   \label{c332}
\end{equation}

\begin{equation}
S_{ch}=S_{ch0}+S_{ch1}\Phi=j_{y}s_{y}\text{.}   \label{c333}
\end{equation}

\bigskip

\begin{itemize}
\item  WEST and EAST boundary:
\end{itemize}

\begin{equation}
S_{ch0}=j_{x}s_{x}\text{,}   \label{c334}
\end{equation}

\begin{equation}
S_{ch1}=0\text{,}   \label{c335}
\end{equation}

\begin{equation}
S_{ch}=S_{ch0}+S_{ch1}\Phi=j_{x}s_{x}\text{.}   \label{c336}
\end{equation}

\bigskip

\bigskip

\paragraph{\texttt{BCPOT=6} ... Prescribe the Value of the Total Current
Flux Density for a Constant Potential (Not Yet Available)}

\bigskip

\bigskip

\paragraph{\texttt{BCPOT=7} ... Prescribe the Potential as a Function of
Other Plasma Parameters (Not Yet Available)}

\bigskip

\bigskip

\paragraph{\texttt{BCPOT=8} ... Prescribe the Value of the Total Electric
Current with a Constant Flux Density}

\bigskip

\begin{equation*}
\mathtt{POTPAR(IB,1)}\equiv J_{x}\text{, }J_{y}
\end{equation*}

\bigskip

\begin{itemize}
\item  SOUTH and NORTH boundary:
\end{itemize}

\begin{equation}
s_{tot}=\sum_{j}^{S\text{ or }N}s_{yj}\text{,}   \label{c337}
\end{equation}

\begin{equation}
S_{ch0}=\left( J_{y}/s_{tot}\right) s_{y}\text{,}   \label{c338}
\end{equation}

\begin{equation}
S_{ch1}=0\text{,}   \label{c339}
\end{equation}

\begin{equation}
S_{ch}=S_{ch0}+S_{ch1}\Phi=\left( J_{y}/s_{tot}\right) s_{y}\text{.} 
\label{c340}
\end{equation}

\bigskip

\begin{itemize}
\item  WEST and EAST boundary:
\end{itemize}

\begin{equation}
s_{tot}=\sum_{j}^{W\text{ or }E}s_{xj}\text{,}   \label{c341}
\end{equation}

\begin{equation}
S_{ch0}=\left( J_{x}/s_{tot}\right) s_{x}\text{,}   \label{c342}
\end{equation}

\begin{equation}
S_{ch1}=0\text{,}   \label{c343}
\end{equation}

\begin{equation}
S_{ch}=S_{ch0}+S_{ch1}\Phi=\left( J_{x}/s_{tot}\right) s_{x}\text{.} 
\label{c344}
\end{equation}

\bigskip

\paragraph{\texttt{BCPOT=9} ... Prescribe the Value of the Decay Length for
the Potential (Only for the SOUTH and NORTH Boundary)}

\bigskip

\begin{equation*}
\mathtt{POTPAR(IB,1)}\equiv l_{\Phi}=\Phi/\partial_{y}\Phi 
\end{equation*}

\bigskip

\begin{equation}
S_{ch0}=0\text{,}   \label{c345}
\end{equation}

\begin{equation}
S_{ch1}=-\frac{\sigma_{a}s_{y}}{l_{\Phi}e_{0}n_{e}}\text{,}^{\left(
\ast\right) }   \label{c346}
\end{equation}

\begin{equation}
S_{ch}=S_{ch0}+S_{ch1}\Phi=-\frac{\sigma_{a}\Phi s_{y}}{l_{\Phi}e_{0}n_{e}}%
\text{.}^{\left( \ast\right) }   \label{c347}
\end{equation}

\bigskip

$^{\left( \ast\right) }$WARNING: PROBABLY\ ERROR\ OR\ THE\ SPECIFIC\
ELECTRIC\ CONDUCTIVITY IS\ GIVEN\ BY $\frac{\sigma_{a}}{e_{0}n_{e}}$!

\bigskip

\bigskip

\subsubsection{Interior Boundary at the ''Dead'' Regions}

\bigskip

\paragraph{Particle Sources}

\bigskip

\begin{equation}
S_{p0\alpha}=\mathcal{V}n_{\alpha\left( is\right) }\text{,}   \label{c348}
\end{equation}

\begin{equation}
S_{p1\alpha}=-\mathcal{V}\text{,}   \label{c349}
\end{equation}

\begin{equation}
S_{p\alpha}=S_{p0\alpha}+S_{p1\alpha}n_{\alpha}=\mathcal{V}\left(
n_{\alpha\left( is\right) }-n_{\alpha}\right) \text{.}   \label{c350}
\end{equation}

\bigskip

\paragraph{Momentum Sources}

\bigskip

\begin{equation}
S_{m0\alpha}=0\text{,}   \label{c351}
\end{equation}

\begin{equation}
S_{m1\alpha}=-\mathcal{VN}_{e}A_{\alpha}s_{y}\text{,}   \label{c352}
\end{equation}

\begin{equation}
S_{m\alpha}=S_{m0\alpha}+S_{m1\alpha}u_{\alpha}=-\mathcal{VN}_{e}A_{\alpha
}s_{y}u_{\alpha}\text{.}   \label{c353}
\end{equation}

\bigskip

\paragraph{Electron Heat Source}

\bigskip

\begin{equation}
S_{he0}=\mathcal{V}n_{e}\widehat{T}_{e\left( is\right) }eV\text{,} 
\label{c354}
\end{equation}

\begin{equation}
S_{he1}=-\mathcal{V}n_{e}\text{,}   \label{c355}
\end{equation}

\begin{equation}
S_{he}=S_{he0}+S_{he1}\widehat{T}_{e}=\mathcal{V}n_{e}\left( \widehat
{T}_{e\left( is\right) }eV-\widehat{T}_{e}\right) \text{.}   \label{c356}
\end{equation}

\bigskip

\paragraph{Ion Heat Source}

\bigskip

\begin{equation}
S_{hi0}=\mathcal{V}n_{e}\widehat{T}_{i\left( is\right) }eV\text{,} 
\label{c357}
\end{equation}

\begin{equation}
S_{hi1}=-\mathcal{V}n_{e}\text{,}   \label{c358}
\end{equation}

\begin{equation}
S_{hi}=S_{hi0}+S_{hi1}\widehat{T}_{i}=\mathcal{V}n_{e}\left( \widehat
{T}_{i\left( is\right) }eV-\widehat{T}_{i}\right) \text{.}   \label{c359}
\end{equation}

\bigskip

\paragraph{Electric Charge Source}

\bigskip

\begin{equation}
S_{ch0}=\mathcal{V}\Phi_{\left( is\right) }\text{,}   \label{c360}
\end{equation}

\begin{equation}
S_{ch1}=-\mathcal{V}\text{,}   \label{c361}
\end{equation}

\begin{equation}
S_{ch}=S_{ch0}+S_{ch1}\Phi=\mathcal{V}\left( \Phi_{\left( is\right)
}-\Phi\right) \text{.}   \label{c362}
\end{equation}

\bigskip

\bigskip

\section{Conclusion}

\bigskip

The existing expressions for the boundary conditions in the B2.5 code cannot
adequately reproduce even relatively simple physical models of the
magnetized boundary plasma. They are based on some unphysical or
insufficiently justified assumptions about the boundary values of the
particle, momentum and energy fluxes as functions of the plasma parameters
calculated by the bulk transport equations. In order to improve the
reliability and accuracy of the B2.5 code, following tasks related to the
boundary conditions should be accomplished:

\begin{itemize}
\item  realistic kinetic or improved fluid models, which take into account
all relevant processes in the boundary plasma, such as solid surface
processes, plasma particle collisions, electric and diamagnetic drifts etc.,
should be developed;

\item  analytical expressions for the boundary conditions, which can be
directly applied into the B2.5 code, should be derived from these improved
models of the boundary plasma;

\item  expressions for the boundary sources should be properly changed, so
that these new analytical expressions for the boundary conditions can be
implemented;

\item  expressions for the boundary sources should be modified, so that
according to the theoretical models some more complicated non-linear
dependences of the boundary conditions on the plasma parameters, can be
implemented.
\end{itemize}

\bigskip

\bigskip

\begin{thebibliography}{99}
\bibitem{pic89}  \label{mpic89}Theilhaber, K., Birdsall, C. K., \textit{%
Phys. Fluids B }\textbf{1 }(1989) 2244, 2260

\bibitem{anders90}  \label{manders90}Anders, A., \textit{A Formulary for
Plasma Physics, }Academie-Verlag Berlin, Berlin, 1990

\bibitem{matt90}  \label{mmatt90}Matthews, G. F., Fielding, S. J.,
McCracken, G. M., Pitcher, C. S., Stangeby, P. C., Ulrickson, M., \textit{%
Plasma Phys. Control. Fusion }\textbf{32 }(1990), 1301

\bibitem{riemann91}  \label{mriemann91}Riemann, K.-U., \textit{J. Phys. D:
Appl. Phys.} \textbf{24 }(1991), 493

\bibitem{stanpitch}  \label{mstanpitch}Stangeby, P. C., Pitcher, C. S.,
Elder, J. D., \textit{Nucl. Fusion }\textbf{32 }(1992), 2079

\bibitem{baelmans}  \label{mbaelmans}Baelmans, M., \textit{Code Improvements
and Applications of a Two-Dimensional Edge Plasma Model for Toroidal
Devices, }PhD Thesis, Universiteit Leuven, Belgium, 1993, published as
Report No. J\"{u}l-2891 and as Report LPP-ERM/KMS No. 100 in 1994

\bibitem{connor}  \label{mconnor}Connor, J. W., Wilson, H. R., \textit{%
Plasma Phys. Control. Fusion }\textbf{36 }(1994), 719

\bibitem{chankin94}  \label{mchankin94}Chankin, A. V., Stangeby, P. C., 
\textit{Plasma Phys. Control. Fusion }\textbf{36 }(1994), 1485

\bibitem{matt94}  \label{mmatt94}Matthews, G. F., \textit{Plasma Phys.
Control. Fusion }\textbf{36 }(1994), 1595

\bibitem{bohm95}  \label{mbohm95}Stangeby, P. C., \textit{Phys. Plasmas }%
\textbf{2 }(1995), 702

\bibitem{stangeby95}  \label{mstangeby95}Stangeby, P. C., Chankin, A. V., 
\textit{Phys. Plasmas }\textbf{2 }(1995), 707

\bibitem{riemann97}  \label{mriemann97}Riemann, K.-U., \textit{Phys. Plasmas 
}\textbf{4 }(1997), 4158

\bibitem{pitch97}  \label{mpitch97}Pitcher, C. S., Stangeby, P. C., \textit{%
Plasma Phys. Control. Fusion }\textbf{39 }(1997), 779

\bibitem{takamura}  \label{mtakamura}Takamura, S., Ye, M. Y., Kuwabara, T.,
Ohno, N., \textit{Phys. Plasmas }\textbf{5 }(1998), 2151

\bibitem{rognlien1}  \label{mrognlien1}Rognlien, T. D., Ryutov, D. D.,
Mattor, N., Porter, G. D., \textit{Phys. Plasmas }\textbf{6 }(1999), 1851

\bibitem{rognlien2}  \label{mrognlien2}Rognlien, T. D., Porter, G. D.,
Ryutov, D. D., \textit{J. Nucl. Mater.} \textbf{266-269} (1999), 654

\bibitem{daube99}  \label{mdaube99}Daube, Th., Riemann, K.-U., \textit{Phys.
Plasmas }\textbf{6 }(1999), 2409

\bibitem{coster}  \label{mcoster}Coster, D. P., \textit{%
b2.parameters.description}, Aug. 3, 1998, IPP Garching, 1999
\end{thebibliography}

\end{document}

--------------22335D01B792A8F762D239EA--


